\chapter{Schemes and Coherent Sheaves}\label{schemes}
\chapterstatus{complete}

\section{Introduction}\label{schemes:section-introduction}

A scheme is a locally ringed space that is locally the spectrum of a ring. Schemes extend the varieties of Chapter~\ref{varieties} in three
directions that Hodge theory needs. They work over arbitrary base rings, so that a variety defined over $\QQ$, or a family over a base, is a single
object; this is how the conjugates $X^\sigma$ of the theory of absolute Hodge classes are defined (Chapter~\ref{absolute-hodge}). They allow nilpotent
functions, which carry infinitesimal information such as tangent vectors and deformations. And they have non-closed points, which make the
generic behaviour of families visible. References: \cite[Chapter II]{hartshorne-ag}, \cite{stacks-project}.

\noindent\textbf{Prerequisites.} Chapter~\ref{varieties} (Algebraic Varieties) and Chapter~\ref{commutative-algebra} (Commutative Algebra).

Rings are commutative with $1$. For a ring $A$ we use the notation $V(I)$, $D(f)$ of Definition~\ref{commutative-algebra:definition-spectrum}.
For a prime $\mathfrak{p}$, $k(\mathfrak{p}) = A_\mathfrak{p}/\mathfrak{p}A_\mathfrak{p}$ is its residue field.

\section{The spectrum of a ring}\label{schemes:section-spectrum}

\begin{proposition}\label{schemes:proposition-spec-topology}
Let $A$ be a ring and $X = \Spec A$.
\begin{enumerate}
\item The sets $V(I)$ are the closed sets of a topology, the \emph{Zariski topology}: $V(IJ) = V(I \cap J) = V(I) \cup V(J)$ and
$V(\sum I_\alpha) = \bigcap V(I_\alpha)$. Moreover $V(I) \subseteq V(J)$ if and only if $\sqrt J \subseteq \sqrt I$.
\item The sets $D(f)$ form a basis of the topology, $D(f) \cap D(g) = D(fg)$, and $D(f) \subseteq D(g)$ if and only if $f \in \sqrt{(g)}$.
\item The closure of $\{\mathfrak{p}\}$ is $V(\mathfrak{p})$. So the closed points are the maximal ideals, and $\mathfrak{p} \mapsto V(\mathfrak{p})$ is a bijection
from $X$ to the set of irreducible closed subsets; $\mathfrak{p}$ is the unique \emph{generic point} of $V(\mathfrak{p})$.
\item $X$ and every $D(f)$ are quasi-compact.
\item A ring map $\varphi \colon A \to B$ induces a continuous map $\Spec B \to \Spec A$, $\mathfrak{q} \mapsto \varphi^{-1}(\mathfrak{q})$, and the preimage
of $D(f)$ is $D(\varphi(f))$.
\end{enumerate}
\end{proposition}

\begin{proof}
(1) A prime contains $IJ$ if and only if it contains $I$ or $J$, and it contains $\sum I_\alpha$ if and only if it contains every $I_\alpha$. The last
statement holds because $\sqrt I$ is the intersection of the primes containing $I$
(Proposition~\ref{commutative-algebra:proposition-nilradical}). (2) $X \setminus V(I) = \bigcup_{f \in I}D(f)$, and $\mathfrak{p}$ omits $fg$ if and
only if it omits both. $D(f) \subseteq D(g)$ means $V(g) \subseteq V(f)$, that is, $f \in \sqrt{(g)}$ by (1). (3) The smallest $V(I)$ containing
$\mathfrak{p}$ is $V(\mathfrak{p})$. If $V(I)$ is irreducible, then $\sqrt I$ is prime: if $fg \in \sqrt I$, then $V(I) = (V(I) \cap V(f)) \cup (V(I) \cap V(g))$,
so $V(I) \subseteq V(f)$, say, and $f \in \sqrt I$. Then $V(I) = V(\sqrt I)$ is the closure of the point $\sqrt I$, which is unique because
$V(\mathfrak{p}) = V(\mathfrak{q})$ forces $\mathfrak{p} = \mathfrak{q}$. (4) If $X = \bigcup_\alpha D(f_\alpha)$, then no prime contains all $f_\alpha$, so they
generate the unit ideal, and $1 = \sum_{i=1}^ra_if_{\alpha_i}$ involves finitely many. For $D(f) \cong \Spec A_f$ (by
Proposition~\ref{commutative-algebra:proposition-localisation}(3)) the same argument applies. (5) $\varphi^{-1}(\mathfrak{q})$ omits $f$ exactly when
$\mathfrak{q}$ omits $\varphi(f)$.
\end{proof}

\begin{definition}\label{schemes:definition-structure-sheaf}
The \emph{structure sheaf} $\mathcal{O} = \mathcal{O}_{\Spec A}$ assigns to an open $U \subseteq \Spec A$ the ring of functions
$s \colon U \to \bigsqcup_{\mathfrak{p} \in U}A_\mathfrak{p}$ with $s(\mathfrak{p}) \in A_\mathfrak{p}$ that are \emph{locally quotients}: every point of $U$ has a
neighbourhood $V \subseteq U$ and $a, f \in A$ with $f \notin \mathfrak{p}$ and $s(\mathfrak{p}) = a/f$ in $A_\mathfrak{p}$ for all $\mathfrak{p} \in V$. The
locally ringed space $(\Spec A, \mathcal{O})$ is the \emph{affine scheme} of $A$.
\end{definition}

\begin{theorem}\label{schemes:theorem-structure-sheaf}
Let $A$ be a ring.
\begin{enumerate}
\item For every $\mathfrak{p}$, the stalk $\mathcal{O}_\mathfrak{p}$ is isomorphic to $A_\mathfrak{p}$; in particular $\Spec A$ is a locally ringed space.
\item For every $f \in A$, $\mathcal{O}(D(f)) \cong A_f$; in particular $\Gamma(\Spec A, \mathcal{O}) \cong A$.
\end{enumerate}
\end{theorem}

\begin{proof}
(1) Evaluation at $\mathfrak{p}$ gives $\mathcal{O}_\mathfrak{p} \to A_\mathfrak{p}$. It is surjective, since $a/f$ with $f \notin \mathfrak{p}$ defines a section on
$D(f)$. It is injective: if $s$, $t$ agree at $\mathfrak{p}$ and are $a/f$, $b/g$ near $\mathfrak{p}$, then $h(ag - bf) = 0$ for some $h \notin \mathfrak{p}$, so
$a/f = b/g$ in every $A_\mathfrak{q}$ with $\mathfrak{q} \in D(fgh)$, a neighbourhood of $\mathfrak{p}$.

(2) Define $\psi \colon A_f \to \mathcal{O}(D(f))$ by $a/f^n \mapsto (\mathfrak{p} \mapsto a/f^n)$. \emph{Injective:} if $a/f^n = 0$ in every $A_\mathfrak{p}$,
$\mathfrak{p} \in D(f)$, then the annihilator of $a$ is contained in no prime of $D(f)$. So $V(\mathrm{Ann}(a)) \subseteq V(f)$, and $f \in \sqrt{\mathrm{Ann}(a)}$,
that is, $f^ma = 0$ and $a/f^n = 0$ in $A_f$. \emph{Surjective:} let $s \in \mathcal{O}(D(f))$. Cover $D(f)$ by open sets on which $s = a_i/g_i$. Shrinking
them to basic open sets $D(h_i) \subseteq D(g_i)$, and writing $h_i^m = cg_i$, so that $a_i/g_i = ca_i/h_i^m$, we may assume $s = a_i/h_i$ on $D(h_i)$. By
(4) of Proposition~\ref{schemes:proposition-spec-topology}, finitely many $D(h_1), \ldots, D(h_r)$ suffice. On $D(h_ih_j)$ both $a_i/h_i$ and $a_j/h_j$
represent $s$. By the injectivity just proved, applied to $D(h_ih_j)$, we get $(h_ih_j)^N(a_ih_j - a_jh_i) = 0$, and we may take one $N$ for all pairs.
Replacing $a_i$ by $a_ih_i^N$ and $h_i$ by $h_i^{N+1}$ we get $a_ih_j = a_jh_i$. Since $D(f) \subseteq \bigcup D(h_i)$, $f^n = \sum_ib_ih_i$ for some $n$ and
$b_i$. Let $a = \sum_ib_ia_i$. Then $h_ja = \sum_ib_ia_ih_j = \sum_ib_ih_ia_j = f^na_j$, so $a/f^n = a_j/h_j$ on $D(h_j)$ for every $j$, and $s = \psi(a/f^n)$.
\end{proof}

\begin{theorem}\label{schemes:theorem-morphisms-to-affine}
For every locally ringed space $X$ and ring $A$, the map
\[
\Hom_{\mathrm{LRS}}(X, \Spec A) \to \Hom_{\mathrm{Ring}}(A, \Gamma(X, \mathcal{O}_X)), \qquad (f, f^\#) \mapsto f^\#_{\Spec A},
\]
is bijective. In particular $A \mapsto \Spec A$ is an anti-equivalence between rings and affine schemes, and a ring map $\varphi$ induces the
morphism $\Spec\varphi$ whose map on stalks is the localisation $A_{\varphi^{-1}(\mathfrak{q})} \to B_\mathfrak{q}$.
\end{theorem}

\begin{proof}
Given $\psi \colon A \to \Gamma(X, \mathcal{O}_X)$, define $f(x) = \psi_x^{-1}(\mathfrak{m}_x)$, where $\psi_x \colon A \to \mathcal{O}_{X,x}$. This is a prime. Then
$f^{-1}(D(a)) = X_{\psi(a)} = \{x : \psi(a)_x \notin \mathfrak{m}_x\}$, which is open, since a germ that is a unit is a unit nearby. So $f$ is continuous. On
$X_{\psi(a)}$ the section $\psi(a)$ is invertible, so $\psi$ induces $A_a \to \mathcal{O}_X(X_{\psi(a)})$. By Theorem~\ref{schemes:theorem-structure-sheaf}
these maps define a morphism of sheaves $\mathcal{O}_{\Spec A} \to f_*\mathcal{O}_X$ on the basis of basic open sets, hence a morphism of sheaves. On the stalk
at $x$ it is the map $A_{f(x)} \to \mathcal{O}_{X,x}$ induced by $\psi_x$, which is local by the definition of $f(x)$. Conversely, let $(f, f^\#)$ be a
morphism with global map $\psi$. Since $f^\#_x \colon A_{f(x)} \to \mathcal{O}_{X,x}$ is local and compatible with $\psi_x$, we get
$f(x)A_{f(x)} = (f^\#_x)^{-1}(\mathfrak{m}_x)$, so $f(x) = \psi_x^{-1}(\mathfrak{m}_x)$. On $D(a)$ the map $f^\#$ is determined by $\psi$ through the universal property
of $A_a$. So the two constructions are inverse. For $X = \Spec B$, $\Gamma(X, \mathcal{O}_X) = B$ by Theorem~\ref{schemes:theorem-structure-sheaf}.
\end{proof}

\section{Schemes and their properties}\label{schemes:section-schemes}

\begin{definition}\label{schemes:definition-scheme}
A \emph{scheme} is a locally ringed space $(X, \mathcal{O}_X)$ in which every point has an open neighbourhood $U$ such that $(U, \mathcal{O}_X|_U)$ is
isomorphic to an affine scheme. A \emph{morphism} of schemes is a morphism of locally ringed spaces. For a scheme $S$, a scheme \emph{over $S$} is a
morphism $X \to S$, and morphisms over $S$ commute with the structure maps. We write $X(T) = \Hom_S(T, X)$ for the \emph{$T$-valued points}, and
$X(R) = X(\Spec R)$. The residue field at $x$ is $k(x) = \mathcal{O}_{X,x}/\mathfrak{m}_x$.
\end{definition}

\begin{proposition}\label{schemes:proposition-gluing}
Let $X_i$ be schemes, $U_{ij} \subseteq X_i$ open subschemes with $U_{ii} = X_i$, and $\varphi_{ij} \colon U_{ij} \to U_{ji}$ isomorphisms such that
$\varphi_{ij}(U_{ij} \cap U_{ik}) = U_{ji} \cap U_{jk}$ and $\varphi_{ik} = \varphi_{jk} \circ \varphi_{ij}$ on $U_{ij} \cap U_{ik}$. Then there is a scheme $X$,
unique up to unique isomorphism, with open subschemes $X_i$ glued along the $\varphi_{ij}$. Morphisms defined on the $X_i$ that agree on overlaps
glue to a morphism $X \to Y$.
\end{proposition}

\begin{proof}
Glue the topological spaces along the homeomorphisms $\varphi_{ij}$, then glue the sheaves, which is possible by the cocycle condition. The result is
locally ringed and locally affine because each $X_i$ is. Morphisms of locally ringed spaces glue because continuous maps and sheaf maps do
(Proposition~\ref{sheaves:proposition-modules}).
\end{proof}

\begin{example}\label{schemes:example-schemes}
\begin{enumerate}
\item $\Spec k$ for a field $k$ is a point. $\Spec\ZZ$ has closed points $(p)$ with residue fields $\mathbb{F}_p$, and a generic point $(0)$ with residue
field $\QQ$, dense in $\Spec\ZZ$.
\item \emph{Affine space} over $A$ is $\Aa^n_A = \Spec A[x_1, \ldots, x_n]$. For $A = k$ algebraically closed, its closed points are the points of
$k^n$ (Theorem~\ref{commutative-algebra:theorem-nullstellensatz}), and it has one further point for each irreducible closed subset of dimension
$> 0$, namely its generic point.
\item \emph{Projective space} $\PP^n_A$ is glued from $U_i = \Spec A[x_0/x_i, \ldots, x_n/x_i]$ along $U_i \supseteq D(x_j/x_i) \cong D(x_i/x_j) \subseteq U_j$. It is
also $\Proj A[x_0, \ldots, x_n]$ (Section~\ref{schemes:section-proj}).
\item The \emph{affine line with doubled origin} is glued from two copies of $\Aa^1_k$ along the identity of $\Aa^1_k \setminus \{0\}$. It is a scheme that is
not isomorphic to any variety (Example~\ref{schemes:example-not-separated}).
\item $\Spec k[\varepsilon]/(\varepsilon^2)$ is a single point with a non-reduced local ring. A morphism $\Spec k[\varepsilon]/(\varepsilon^2) \to X$ over $k$ is a
$k$-point $x$ of $X$ together with a tangent vector in $(\mathfrak{m}_x/\mathfrak{m}_x^2)^\vee$ (Exercise~\ref{schemes:exercise-dual-numbers}).
\end{enumerate}
\end{example}

\begin{definition}\label{schemes:definition-properties}
A scheme $X$ is \emph{reduced} if all local rings are reduced, \emph{irreducible} if its space is, \emph{integral} if every $\mathcal{O}_X(U)$, $U \neq \emptyset$,
is a domain, \emph{locally Noetherian} if it has an affine open cover by spectra of Noetherian rings, and \emph{Noetherian} if moreover it is
quasi-compact. A morphism $f \colon X \to Y$ is \emph{locally of finite type} if there are affine open covers $Y = \bigcup\Spec B_i$ and
$f^{-1}(\Spec B_i) = \bigcup_j\Spec A_{ij}$ with each $A_{ij}$ a finitely generated $B_i$-algebra, \emph{of finite type} if moreover finitely many $A_{ij}$
suffice for each $i$, and \emph{finite} if $f^{-1}(\Spec B_i) = \Spec A_i$ with $A_i$ a finite $B_i$-module. An \emph{open immersion} is an isomorphism onto
an open subscheme. A \emph{closed immersion} is a morphism that is a homeomorphism onto a closed subset and such that
$f^\# \colon \mathcal{O}_Y \to f_*\mathcal{O}_X$ is surjective. A \emph{closed subscheme} is an equivalence class of closed immersions.
\end{definition}

\begin{proposition}\label{schemes:proposition-integral}
A scheme is integral if and only if it is reduced and irreducible. An integral scheme $X$ has a unique generic point $\eta$, and its
\emph{function field} is $K(X) = \mathcal{O}_{X,\eta}$, which is the fraction field of $A$ for every nonempty affine open $\Spec A \subseteq X$.
\end{proposition}

\begin{proof}
If $X$ is integral, it is reduced. If it were reducible, there would be disjoint nonempty open sets $U_1$, $U_2$, and
$\mathcal{O}(U_1 \sqcup U_2) = \mathcal{O}(U_1) \times \mathcal{O}(U_2)$ would not be a domain. Conversely, let $X$ be reduced and irreducible, and let
$U = \Spec A$ be a nonempty affine open. Then $A$ is reduced with irreducible spectrum, so its nilradical $0$ is prime
(Proposition~\ref{schemes:proposition-spec-topology}(3)), and $A$ is a domain. For general $U$ and $fg = 0$ in $\mathcal{O}(U)$, $U$ is the union of the
closed sets where the germs of $f$ and of $g$ lie in the maximal ideal. On each affine open these are $V(f)$ and $V(g)$. By irreducibility one of
them, say $V(f)$, is all of $U$, so $f$ is nilpotent on each affine open, hence $0$. The generic point of $X$ is the generic point of each nonempty
open subset, which for $\Spec A$ is the prime $(0)$, with local ring $A_{(0)} = \mathrm{Frac}(A)$.
\end{proof}

\begin{proposition}\label{schemes:proposition-points-over-field}
Let $X$ be a scheme and $K$ a field. Morphisms $\Spec K \to X$ correspond bijectively to pairs of a point $x \in X$ and a field homomorphism
$k(x) \to K$. If $X$ is of finite type over an algebraically closed field $k$, then the closed points of $X$ are exactly the points with $k(x) = k$, and
$X(k)$ is the set of closed points.
\end{proposition}

\begin{proof}
A morphism from the one-point space $\Spec K$ has image a point $x$, and the stalk map $\mathcal{O}_{X,x} \to K$ is local, so it kills $\mathfrak{m}_x$ and
factors through $k(x)$. Conversely such data define a morphism, by Theorem~\ref{schemes:theorem-morphisms-to-affine} applied on an affine
neighbourhood. For $X$ of finite type over $k$, a point $x$ in an affine open $\Spec A$ is closed in $\Spec A$ if and only if its prime is maximal,
if and only if $k(x) = k$, by Zariski's lemma (Theorem~\ref{commutative-algebra:theorem-zariski-lemma}) and algebraic closedness. If $k(x) = k$,
then $x$ is closed in every affine open containing it, and its closure does not meet any open set not containing $x$, since open sets are stable
under generisation. So $x$ is closed in $X$. Conversely, a closed point of $X$ is closed in an affine open containing it.
\end{proof}

\section{Fibre products}\label{schemes:section-fibre-products}

\begin{theorem}\label{schemes:theorem-fibre-products}
Fibre products exist in the category of schemes. For affine schemes,
$\Spec A \times_{\Spec R}\Spec B = \Spec(A \otimes_RB)$.
\end{theorem}

\begin{proof}[Sketch of proof]
For affine schemes this follows from Theorem~\ref{schemes:theorem-morphisms-to-affine} and the universal property of $A \otimes_RB$
(Proposition~\ref{rings:proposition-tensor-algebras}): a morphism $T \to \Spec(A \otimes_RB)$ is a ring map $A \otimes_RB \to \Gamma(T, \mathcal{O}_T)$, that is,
a compatible pair of ring maps from $A$ and $B$. In general one glues. If $X \times_SY$ exists and $U \subseteq X$ is open, then the preimage of $U$ in
$X \times_SY$ is $U \times_SY$. Hence, covering $S$, then $X$ and $Y$ by affine opens, the affine fibre products glue along these canonical open
subschemes by Proposition~\ref{schemes:proposition-gluing}. The omitted verifications are the cocycle conditions, which follow from the uniqueness
of fibre products, and the universal property of the glued scheme, which is local \cite[Theorem II.3.3]{hartshorne-ag}.
\end{proof}

\begin{definition}\label{schemes:definition-fibre}
For a morphism $f \colon X \to Y$ and $y \in Y$, the \emph{fibre} is $X_y = X \times_Y\Spec k(y)$, a scheme over the residue field $k(y)$ whose underlying space
is homeomorphic to $f^{-1}(y)$ \cite[Exercise II.3.10]{hartshorne-ag}. For $Y' \to Y$, the \emph{base change} of $X$ is $X_{Y'} = X \times_YY'$. A property of
morphisms is \emph{stable under base change} if $X_{Y'} \to Y'$ has it whenever $X \to Y$ does.
\end{definition}

\begin{example}\label{schemes:example-fibre-products}
\begin{enumerate}
\item $\Aa^m_R \times_R \Aa^n_R = \Aa^{m+n}_R$, since $R[x] \otimes_RR[y] = R[x, y]$.
\item $\Spec\CC \times_{\Spec\RR}\Spec\CC = \Spec(\CC \otimes_\RR\CC) = \Spec(\CC \times \CC)$, two points. So a fibre product of one-point schemes can have
several points, and the underlying set of a fibre product is not the fibre product of the underlying sets.
\item For $f \colon \Spec\ZZ[i] \to \Spec\ZZ$, the fibre over $(p)$ is $\Spec\mathbb{F}_p[x]/(x^2 + 1)$. It is two reduced points if $p \equiv 1 \bmod 4$, one
point with residue field $\mathbb{F}_{p^2}$ if $p \equiv 3 \bmod 4$, and the non-reduced point $\Spec\mathbb{F}_2[x]/(x + 1)^2$ if $p = 2$. The
generic fibre is $\Spec\QQ(i)$.
\end{enumerate}
\end{example}

\section{Separated and proper morphisms}\label{schemes:section-separated-proper}

\begin{definition}\label{schemes:definition-separated}
A morphism $f \colon X \to Y$ is \emph{separated} if the diagonal $\Delta \colon X \to X \times_YX$ is a closed immersion. It is \emph{proper} if it is
separated, of finite type, and \emph{universally closed}: every base change $X_{Y'} \to Y'$ is a closed map. A scheme over a field $k$ is \emph{complete} if
it is proper over $\Spec k$.
\end{definition}

\begin{proposition}\label{schemes:proposition-affine-separated}
Every morphism of affine schemes is separated. Open and closed immersions are separated, and so are compositions and base changes of separated
morphisms. A morphism $X \to Y$ is separated if and only if the diagonal has closed image.
\end{proposition}

\begin{proof}
For $\Spec A \to \Spec R$ the diagonal corresponds to the multiplication $A \otimes_RA \to A$, which is surjective, so $\Delta$ is a closed immersion (a
surjection of rings $B \to B/I$ gives a homeomorphism onto $V(I)$ and a surjection of sheaves, by Theorem~\ref{schemes:theorem-structure-sheaf}).
For the remaining statements, see \cite[Corollary II.4.2 and Corollary II.4.6]{hartshorne-ag}. The last one holds because $\Delta$ is always a
locally closed immersion: over an affine open of $Y$ and affine opens $U$ of $X$, the diagonal maps $U$ by a closed immersion into the open
subscheme $U \times_YU$.
\end{proof}

\begin{example}\label{schemes:example-not-separated}
The affine line with doubled origin $X$ is not separated over $k$. Let $i_1, i_2 \colon \Aa^1_k \to X$ be the two open immersions. The morphism
$(i_1, i_2) \colon \Aa^1_k \to X \times_kX$ maps $t \neq 0$ into the diagonal but maps $0$ to the pair of the two origins, which is not in the diagonal. So
the preimage of $\Delta(X)$ is $\Aa^1_k \setminus \{0\}$, which is not closed. Hence $\Delta(X)$ is not closed.
\end{example}

\begin{theorem}[Valuative criteria]\label{schemes:theorem-valuative}
Let $f \colon X \to Y$ be a morphism of finite type with $X$ Noetherian. For every valuation ring $R$ with fraction field $K$ and every
commutative square formed by $\Spec K \to X$ and $\Spec R \to Y$, consider lifts $\Spec R \to X$ making both triangles commute. Then $f$ is separated
if and only if there is at most one lift for every such square, and proper if and only if there is exactly one. It suffices to consider discrete
valuation rings.
\end{theorem}

This is \cite[Theorems II.4.3 and II.4.7 and Exercise II.4.11]{hartshorne-ag}. Geometrically: a punctured disc $\Spec K$ in $X$ whose image extends
over the centre of $\Spec R$ in $Y$ extends in at most one way (separated: limits are unique) or in exactly one way (proper: limits exist).

\begin{theorem}\label{schemes:theorem-projective-proper}
For every ring $A$, $\PP^n_A \to \Spec A$ is proper. Consequently every \emph{projective} morphism, that is, a closed immersion into some $\PP^n_Y$
followed by the projection, is proper, and every projective variety is complete.
\end{theorem}

\begin{proof}
Since $\PP^n_A = \PP^n_\ZZ \times \Spec A$ and properness is stable under base change \cite[Corollary II.4.8]{hartshorne-ag}, we may take $A = \ZZ$. The
morphism is of finite type, being covered by the $n + 1$ affine spaces $U_i$. It is separated: on $U_i \times U_j$ the diagonal is
$U_i \cap U_j \to U_i \times U_j$, given by the ring map $\ZZ[x/x_i] \otimes \ZZ[x/x_j] \to \ZZ[x/x_i]_{x_j/x_i}$, which is surjective because $x_i/x_j$ is the
image of an element of the second factor. By Proposition~\ref{schemes:proposition-affine-separated} this suffices. For properness, use
Theorem~\ref{schemes:theorem-valuative}. A morphism $\Spec K \to \PP^n_\ZZ$ lands in some $U_i$, so by
Theorem~\ref{schemes:theorem-morphisms-to-affine} it is given by elements $a_j = x_j/x_i \in K$ with $a_i = 1$. Let $v$ be the valuation of $R$
(with $v(0) = \infty$) and choose $m$ with $v(a_m)$ minimal; then $a_m \neq 0$. All $b_j = a_j/a_m$ lie in $R$, and the ring map
$\ZZ[x_0/x_m, \ldots, x_n/x_m] \to R$, $x_j/x_m \mapsto b_j$, defines $\Spec R \to U_m$. On $\Spec K$ it agrees with the given morphism, because on
$U_i \cap U_m$ the coordinates $x_j/x_i = (x_j/x_m)(x_m/x_i)$ are sent to $b_j/b_i = a_j$. Uniqueness follows from separatedness. Closed immersions are proper, and compositions of proper morphisms are proper \cite[Corollary II.4.8]{hartshorne-ag}.
\end{proof}

\section{Projective schemes}\label{schemes:section-proj}

\begin{definition}\label{schemes:definition-proj}
Let $S = \bigoplus_{d \geq 0}S_d$ be a graded ring and $S_+ = \bigoplus_{d > 0}S_d$. The set $\Proj S$ consists of the homogeneous primes not containing
$S_+$, with the topology whose closed sets are the $V_+(\mathfrak{a})$, $\mathfrak{a}$ homogeneous. For homogeneous $f \in S_+$ let
$D_+(f) = \{\mathfrak{p} : f \notin \mathfrak{p}\}$, and let $S_{(f)}$ be the degree-zero part of $S_f$, and $S_{(\mathfrak{p})}$ the degree-zero part of the
localisation at the homogeneous elements not in $\mathfrak{p}$. The structure sheaf consists of the functions $s(\mathfrak{p}) \in S_{(\mathfrak{p})}$ that are
locally quotients $a/f$ with $a$, $f$ homogeneous of the same degree.
\end{definition}

\begin{proposition}\label{schemes:proposition-proj}
$\Proj S$ is a scheme. The stalk at $\mathfrak{p}$ is $S_{(\mathfrak{p})}$, the sets $D_+(f)$ form an open cover, and $D_+(f) \cong \Spec S_{(f)}$. There is a
morphism $\Proj S \to \Spec S_0$. If $S = A[x_0, \ldots, x_n]$, then $\Proj S = \PP^n_A$, and if $S = A[x_0, \ldots, x_n]/I$ with $I$ homogeneous, then
$\Proj S$ is a closed subscheme of $\PP^n_A$.
\end{proposition}

\begin{proof}
The stalks are computed as in Theorem~\ref{schemes:theorem-structure-sheaf}(1). Every $\mathfrak{p} \in \Proj S$ omits some homogeneous element of
$S_+$, so the $D_+(f)$ cover. Let $f$ have degree $e > 0$. The map $\mathfrak{p} \mapsto (\mathfrak{p}S_f) \cap S_{(f)}$ is a homeomorphism
$D_+(f) \to \Spec S_{(f)}$. Its inverse sends $\mathfrak{q}$ to the homogeneous ideal whose degree $d$ part is $\{a \in S_d : a^e/f^d \in \mathfrak{q}\}$; checking
that this is a homogeneous prime and that the two maps are inverse homeomorphisms is routine \cite[Proposition II.2.5]{hartshorne-ag}. Under it,
$S_{(\mathfrak{p})}$ corresponds to the localisation of $S_{(f)}$ at $\mathfrak{q}$, and the locally-quotient descriptions of the two structure sheaves
match, so the ringed spaces are isomorphic. The inclusions $S_0 \to S_{(f)}$ glue to the morphism to $\Spec S_0$. For $S = A[x_0, \ldots, x_n]$,
$S_{(x_i)} = A[x_0/x_i, \ldots, x_n/x_i]$, and the gluing agrees with Example~\ref{schemes:example-schemes}(3). A surjection $S \to S/I$ induces
surjections $S_{(x_i)} \to (S/I)_{(x_i)}$, hence closed immersions $\Spec(S/I)_{(x_i)} \to \Spec S_{(x_i)}$, which glue to a closed immersion
$\Proj S/I \to \PP^n_A$.
\end{proof}

\section{Varieties as schemes}\label{schemes:section-varieties}

\begin{theorem}\label{schemes:theorem-varieties-schemes}
Let $k$ be algebraically closed. There is a fully faithful functor $t$ from the category of varieties over $k$ (Chapter~\ref{varieties}) to the category of
schemes over $k$. The space $t(V)$ is the set of irreducible closed subsets of $V$, and $V$ is homeomorphic to the subspace of closed points of $t(V)$.
The image consists of the integral separated schemes of finite type over $k$ that admit an immersion into some $\PP^n_k$, and $t(V)$ is proper over $k$ if
and only if $V$ is projective.
\end{theorem}

\begin{proof}[Sketch of proof]
For an affine variety $V$ with coordinate ring $A$, $t(V) = \Spec A$. Irreducible closed subsets of $V$ correspond to primes of $A$
(Proposition~\ref{varieties:proposition-irreducible-components}), and points to maximal ideals. By Theorem~\ref{varieties:theorem-regular-functions-affine}
and Theorem~\ref{schemes:theorem-structure-sheaf}, the regular functions on an open subset $U \subseteq V$ are the sections of $\mathcal{O}_{\Spec A}$ over the
corresponding open subset. In general $t(V)$ is glued from affine pieces. Full faithfulness reduces, by gluing, to the affine case, where both
categories are anti-equivalent to the category of finitely generated domains (Proposition~\ref{varieties:proposition-morphisms-to-affine} and
Theorem~\ref{schemes:theorem-morphisms-to-affine}). The omitted verifications are that the construction is compatible with gluing and that
quasi-projective varieties give separated schemes, which is Proposition~\ref{schemes:proposition-affine-separated} and
Theorem~\ref{schemes:theorem-projective-proper}; see \cite[Proposition II.2.6 and Proposition II.4.10]{hartshorne-ag}.
\end{proof}

\begin{remark}\label{schemes:remark-variety-convention}
From now on, a \emph{variety over a field $k$} is an integral separated scheme of finite type over $k$. For $k$ algebraically closed and quasi-projective
varieties this agrees with Chapter~\ref{varieties} by Theorem~\ref{schemes:theorem-varieties-schemes}. There are complete varieties that are not
projective (Nagata, Hironaka), but every complete curve is projective, and so is every smooth complete surface; see \cite[Appendix B]{hartshorne-ag}.
\end{remark}

\section{Flat and smooth morphisms}\label{schemes:section-flat-smooth}

\begin{definition}\label{schemes:definition-flat}
A morphism $f \colon X \to Y$ is \emph{flat} at $x$ if $\mathcal{O}_{X,x}$ is a flat $\mathcal{O}_{Y,f(x)}$-module (Definition~\ref{rings:definition-flat}),
and \emph{flat} if it is flat at every point. It is \emph{faithfully flat} if moreover it is surjective.
\end{definition}

\begin{proposition}\label{schemes:proposition-flat}
\begin{enumerate}
\item Open immersions are flat, and $\Aa^n_Y \to Y$ and $\PP^n_Y \to Y$ are flat. Compositions and base changes of flat morphisms are flat.
\item Let $R$ be a principal ideal domain, for instance a discrete valuation ring. An $R$-module is flat if and only if it is torsion-free. Hence
$X \to \Spec R$ is flat if and only if every local ring $\mathcal{O}_{X,x}$ is torsion-free over $R$. If $X$ is integral, this holds if and only if
$X$ dominates $\Spec R$.
\end{enumerate}
\end{proposition}

\begin{proof}
(1) Localisations are flat (Proposition~\ref{commutative-algebra:proposition-localisation}(2)), and so are free modules such as $A[x_1, \ldots, x_n]$ over $A$;
local rings of $\Aa^n_A$ are localisations of $A[x]$. The statements on composition and base change are \cite[Proposition III.9.2]{hartshorne-ag}.
(2) If $M$ is flat and $0 \neq t \in R$, tensoring the injection $t \colon R \to R$ with $M$ shows that $t$ is injective on $M$. Conversely, a torsion-free
$M$ is the union of its finitely generated submodules, which are torsion-free, hence free (Theorem~\ref{rings:theorem-structure-pid}). A union of flat
submodules is flat, since an element of the kernel of $N' \otimes M \to N \otimes M$ comes from some finitely generated submodule of $M$. For integral $X$,
torsion-freeness fails exactly when some nonzero $t \in R$ vanishes on $X$, that is, when $X$ maps into a closed point.
\end{proof}

\begin{example}\label{schemes:example-flat}
\begin{enumerate}
\item The family $X = \Spec k[t, x, y]/(xy - t) \to \Aa^1_k = \Spec k[t]$ is flat, since $X$ is integral and dominates $\Aa^1$. Its fibres are the smooth conics
$xy = t$ for $t \neq 0$, which degenerate to the pair of lines $xy = 0$.
\item The blow-up $\mathrm{Bl}_0\Aa^2 \to \Aa^2$ is not flat. Its fibres have dimension $0$ except over the origin, where the fibre is $\PP^1$, while flat morphisms
of finite type between varieties have equidimensional fibres \cite[Proposition III.9.5]{hartshorne-ag}.
\item A closed immersion $\Spec A/I \to \Spec A$ with $0 \neq I$ and $A$ a domain is not flat, since $A/I$ is torsion.
\end{enumerate}
\end{example}

\begin{definition}\label{schemes:definition-smooth}
A morphism $f \colon X \to Y$ is \emph{smooth of relative dimension $n$} if every point of $X$ has affine open neighbourhoods $\Spec B$ of $x$ and $\Spec A$ of
$f(x)$, with $f(\Spec B) \subseteq \Spec A$, such that
\[
B \cong A[x_1, \ldots, x_{n+r}]/(g_1, \ldots, g_r) \qquad \text{with } \det\left(\partial g_i/\partial x_{n+j}\right)_{1 \leq i,j \leq r} \text{ invertible in } B.
\]
It is \emph{\'etale} if it is smooth of relative dimension $0$.
\end{definition}

\begin{proposition}\label{schemes:proposition-smooth-nonsingular}
Let $k$ be algebraically closed and $X$ a variety over $k$ of dimension $n$. Then $X \to \Spec k$ is smooth (of relative dimension $n$) if and only if $X$
is nonsingular in the sense of Definition~\ref{varieties:definition-tangent-space}.
\end{proposition}

\begin{proof}
Both conditions are local, so let $X$ be affine. Suppose $B = k[x_1, \ldots, x_{n+r}]/(g_1, \ldots, g_r)$ with the minor invertible, and let $P$ be a closed point.
Every irreducible component of $Z(g_1, \ldots, g_r) \subseteq \Aa^{n+r}$ has dimension $\geq n$ by Theorem~\ref{varieties:theorem-hypersurface-sections} applied $r$
times, so $\dim\mathcal{O}_{X,P} \geq n$. The Jacobian of $(g_i)$ has rank $r$ at $P$, so $\dim T_PX \leq n$ by the computation in
Theorem~\ref{varieties:theorem-jacobian}, which only used generators of the ideal defining $B$. By Lemma~\ref{varieties:lemma-tangent-dimension},
$\mathcal{O}_{X,P}$ is regular of dimension $n$. Conversely, let $X \subseteq \Aa^N$ be nonsingular at $P$ with $I(X) = (f_1, \ldots, f_t)$. The Jacobian has rank
$r = N - n$ at $P$, so after renumbering, $\det(\partial f_i/\partial x_{n+j})_{i,j \leq r}$ is nonzero at $P$, hence invertible on a neighbourhood $D(h)$. The
scheme $X' = Z(f_1, \ldots, f_r) \cap D(h)$ is smooth of relative dimension $n$ by construction. So by the first part its local ring at $P$ is regular of dimension
$n$, in particular a domain (Remark~\ref{commutative-algebra:remark-regular-domain}). Hence $X'$ has a unique irreducible component through $P$, of dimension $n$, which contains
$X \cap D(h)$ and so equals it near $P$. As $X'$ is reduced near $P$, $X = X'$ near $P$.
\end{proof}

\begin{counterexample}\label{schemes:counterexample-regular-not-smooth}
Over a non-perfect field, regular does not imply smooth. Let $p$ be an odd prime, $k = \mathbb{F}_p(t)$ and $X = \Spec k[x, y]/(y^2 - x^p + t)$. At the closed point
$\mathfrak{m} = (y, x^p - t)$, whose residue field is $k(t^{1/p})$, the relation $x^p - t = y^2 \in \mathfrak{m}^2$ shows that $\mathfrak{m}/\mathfrak{m}^2$ is spanned by $y$. So
$X$ is regular there, of dimension $1$. But $X$ is not smooth over $k$. Smoothness is stable under base change, since a presentation as in
Definition~\ref{schemes:definition-smooth} base changes to one. After base change to $\bar k$, with $s = t^{1/p}$, the curve becomes
$y^2 = (x - s)^p$, whose Jacobian $(-p(x - s)^{p-1}, 2y) = (0, 2y)$ vanishes at the point $(s, 0)$. That point is singular by
Theorem~\ref{varieties:theorem-jacobian}, so $X_{\bar k}$ is not smooth by Proposition~\ref{schemes:proposition-smooth-nonsingular}.
\end{counterexample}

\section{Base change and varieties over non-closed fields}\label{schemes:section-base-change}

\begin{definition}\label{schemes:definition-geometric}
Let $X$ be a scheme over a field $k$ and $K \supseteq k$ a field extension. The \emph{base change} is $X_K = X \times_{\Spec k}\Spec K$. The scheme $X$ is
\emph{geometrically reduced}, \emph{irreducible} or \emph{integral} if $X_{\bar k}$ is reduced, irreducible or integral for an algebraic closure $\bar k$. A
\emph{geometric point} is a morphism $\Spec\Omega \to X$ with $\Omega$ algebraically closed.
\end{definition}

\begin{example}\label{schemes:example-non-closed-fields}
\begin{enumerate}
\item $X = \Spec\RR[x]/(x^2 + 1) = \Spec\CC$ is integral over $\RR$, but $X_\CC = \Spec\CC \sqcup \Spec\CC$ is not irreducible.
\item $X = \Spec\mathbb{F}_p(t)[x]/(x^p - t)$ is a point with residue field $\mathbb{F}_p(t^{1/p})$, integral, but $X_{\bar k}$ is non-reduced, since $x^p - t = (x - t^{1/p})^p$.
\item The conic $X = \Spec\RR[x, y]/(x^2 + y^2 + 1)$ has no $\RR$-points, but it is a nonempty integral scheme whose closed points have residue field $\CC$.
$X(\CC)$ is an affine conic with an action of complex conjugation without fixed points.
\item For a variety $X$ over $k$ and a Galois extension $K/k$, $\mathrm{Gal}(K/k)$ acts on $X(K) = \Hom_k(\Spec K, X)$ through its action on $\Spec K$, and
$X(k) \subseteq X(K)^{\mathrm{Gal}(K/k)}$. For a variety over $\QQ$, the set $X(\bar{\QQ})$ with its Galois action is the starting point of the $\ell$-adic
cohomology of Chapter~\ref{etale-cohomology}.
\end{enumerate}
\end{example}

\begin{remark}\label{schemes:remark-spreading-out}
A variety $X$ over $\CC$ is defined by finitely many polynomials, so it is defined over a finitely generated subfield $k_0 \subseteq \CC$: $X = X_0 \times_{k_0}\CC$.
Moreover, $X_0$ extends to a scheme of finite type over a finitely generated $\ZZ$-algebra, a technique called \emph{spreading out}. For an automorphism
$\sigma$ of $\CC$, the \emph{conjugate variety} is $X^\sigma = X \times_{\CC,\sigma}\CC$. It is a scheme isomorphic to $X$ as an abstract scheme but not
over $\CC$, and $X^\sigma(\CC)$ is in general not homeomorphic to $X(\CC)$ \cite{serre-1964}. These constructions are used in the theory of absolute Hodge
classes (Chapter~\ref{absolute-hodge}).
\end{remark}

\section{Quasi-coherent and Coherent Sheaves}\label{schemes:section-coherent-sheaves}
Quasi-coherent sheaves are to schemes what modules are to rings. On an affine scheme $\Spec A$ they are exactly the sheaves $\widetilde M$ attached
to $A$-modules, and in general they are glued from such. Coherent sheaves, the finitely generated ones, include vector bundles, ideal sheaves of
subvarieties, structure sheaves of closed subschemes and sheaves of differentials. Their cohomology is the subject of
Section~\ref{schemes:section-coherent-cohomology}. The following sections develop the basic theory: the equivalence between modules and quasi-coherent sheaves on affine
schemes, closed subschemes and ideal sheaves, the twisting sheaves $\mathcal{O}(n)$ and Serre's theorem that $\mathcal{F}(n)$ is globally generated
for $n \gg 0$, and sheaves of differentials, with the Euler sequence and the canonical sheaf of $\PP^n$. References:
\cite[Sections II.5 and II.8]{hartshorne-ag}, \cite{stacks-project}.

\noindent\textbf{Prerequisites.} Chapter~\ref{schemes} (Schemes).

For a ringed space $(X, \mathcal{O}_X)$, $\mathcal{O}_X$-modules, their tensor products, $\mathcal{H}om$, direct and inverse images are as in
Definition~\ref{sheaves:definition-module-operations}. For $f \in \Gamma(X, \mathcal{O}_X)$ we write $X_f = \{x \in X : f_x \notin \mathfrak{m}_x\}$, an open
set on which $f$ is invertible.

\section{The sheaf associated to a module}\label{schemes:section-tilde}

\begin{definition}\label{schemes:definition-tilde}
Let $A$ be a ring and $M$ an $A$-module. The sheaf $\widetilde M$ on $\Spec A$ assigns to an open $U$ the set of functions
$s \colon U \to \bigsqcup_{\mathfrak{p} \in U}M_\mathfrak{p}$ with $s(\mathfrak{p}) \in M_\mathfrak{p}$ that are locally of the form $m/f$ with $m \in M$ and $f \in A$
nonvanishing on the neighbourhood. It is an $\mathcal{O}_{\Spec A}$-module.
\end{definition}

\begin{proposition}\label{schemes:proposition-tilde}
Let $A$ be a ring, $M$, $N$ modules and $X = \Spec A$.
\begin{enumerate}
\item $(\widetilde M)_\mathfrak{p} \cong M_\mathfrak{p}$ and $\widetilde M(D(f)) \cong M_f$; in particular $\Gamma(X, \widetilde M) = M$.
\item $M \mapsto \widetilde M$ is an exact, fully faithful functor from $A$-modules to $\mathcal{O}_X$-modules, commuting with direct sums, and
$(M \otimes_AN)^\sim \cong \widetilde M \otimes_{\mathcal{O}_X}\widetilde N$.
\item For a ring map $A \to B$ and $f \colon \Spec B \to \Spec A$: $f_*(\widetilde N) \cong ({}_AN)^\sim$ for a $B$-module $N$, and
$f^*(\widetilde M) \cong (B \otimes_AM)^\sim$.
\end{enumerate}
\end{proposition}

\begin{proof}
(1) The proof of Theorem~\ref{schemes:theorem-structure-sheaf} applies word for word with $A$ replaced by $M$ in the numerators. (2) Exactness can be
checked on stalks, where it is exactness of localisation (Proposition~\ref{commutative-algebra:proposition-localisation}(2)). Full faithfulness: a
morphism $\widetilde M \to \widetilde N$ is determined by its effect on global sections, by (1) on basic open sets, and every $A$-linear map induces one.
Direct sums and tensor products commute with localisation, and the sheaf tensor product has stalks $M_\mathfrak{p} \otimes_{A_\mathfrak{p}}N_\mathfrak{p}$; the
natural map is therefore an isomorphism on stalks. (3) For $g \in A$ with image $\bar g \in B$,
$f_*(\widetilde N)(D(g)) = \widetilde N(D(\bar g)) = N_{\bar g} = (N_A)_g$, compatibly with restrictions. For the pullback, both sides represent the functor
$\mathcal{G} \mapsto \Hom_{\mathcal{O}_{\Spec A}}(\widetilde M, f_*\mathcal{G})$ on quasi-coherent $\mathcal{G}$, by adjunction and (2); see
\cite[Proposition II.5.2]{hartshorne-ag}.
\end{proof}

\section{Quasi-coherent sheaves}\label{schemes:section-quasi-coherent}

\begin{definition}\label{schemes:definition-quasi-coherent}
An $\mathcal{O}_X$-module $\mathcal{F}$ on a scheme $X$ is \emph{quasi-coherent} if $X$ has an affine open cover $X = \bigcup\Spec A_i$ with
$\mathcal{F}|_{\Spec A_i} \cong \widetilde{M_i}$ for $A_i$-modules $M_i$. If $X$ is Noetherian, $\mathcal{F}$ is \emph{coherent} if moreover the $M_i$ can be
taken finitely generated.
\end{definition}

\begin{lemma}\label{schemes:lemma-extend-sections}
Let $X$ be a scheme that is Noetherian, or quasi-compact and separated. Let $\mathcal{L}$ be an invertible sheaf, that is, locally isomorphic to $\mathcal{O}_X$, $f \in \Gamma(X, \mathcal{L})$, and
$X_f = \{x : f_x \notin \mathfrak{m}_x\mathcal{L}_x\}$. Let $\mathcal{F}$ be quasi-coherent.
\begin{enumerate}
\item If $s \in \Gamma(X, \mathcal{F})$ restricts to $0$ on $X_f$, then $s \otimes f^n = 0$ in $\Gamma(X, \mathcal{F} \otimes \mathcal{L}^n)$ for some $n$.
\item If $t \in \mathcal{F}(X_f)$, then $t \otimes f^n$ extends to a global section of $\mathcal{F} \otimes \mathcal{L}^n$ for some $n$.
\end{enumerate}
\end{lemma}

\begin{proof}
Cover $X$ by finitely many affine opens $U_i = \Spec A_i$ with $\mathcal{L}|_{U_i} \cong \mathcal{O}_{U_i}$ and $\mathcal{F}|_{U_i} = \widetilde{M_i}$. Every
$U_i \cap U_j$ is quasi-compact: every open subset of a Noetherian scheme is, and in the separated case $U_i \cap U_j$ is affine, being the preimage of
the closed subscheme $\Delta(X)$ under $U_i \times U_j \to X \times X$. Under the trivialisation $f$ becomes $f_i \in A_i$ and
$X_f \cap U_i = D(f_i)$. (1) $s|_{U_i} \in M_i$ vanishes in $(M_i)_{f_i}$ by Proposition~\ref{schemes:proposition-tilde}(1), so $f_i^ns|_{U_i} = 0$ for
large $n$, which we take uniform over the finitely many $i$. (2) $t|_{D(f_i)} \in (M_i)_{f_i}$, so $f_i^nt$ extends to a section $t_i$ over $U_i$. On
$U_i \cap U_j$, which is quasi-compact and so covered by finitely many affine opens, $t_i - t_j$ vanishes on $X_f$. By (1) applied to $U_i \cap U_j$,
$f^m(t_i - t_j) = 0$ there. So the sections $f^mt_i$ glue to a global section of $\mathcal{F} \otimes \mathcal{L}^{n+m}$ extending $t \otimes f^{n+m}$.
\end{proof}

\begin{theorem}\label{schemes:theorem-affine}
Let $X = \Spec A$. An $\mathcal{O}_X$-module $\mathcal{F}$ is quasi-coherent if and only if the natural map $\Gamma(X, \mathcal{F})^\sim \to \mathcal{F}$ is an
isomorphism. Hence $M \mapsto \widetilde M$ is an equivalence between $A$-modules and quasi-coherent $\mathcal{O}_X$-modules, with inverse $\Gamma(X, -)$. A
sequence of quasi-coherent sheaves on $X$ is exact if and only if the sequence of global sections is exact.
\end{theorem}

\begin{proof}
The natural map is defined on $D(g)$ by $M_g \to \mathcal{F}(D(g))$, $m/g^n \mapsto g^{-n}m|_{D(g)}$, where $M = \Gamma(X, \mathcal{F})$. If $\mathcal{F}$ is
quasi-coherent, apply Lemma~\ref{schemes:lemma-extend-sections} with $\mathcal{L} = \mathcal{O}_X$ and $f = g$ to the affine, hence quasi-compact
and separated, scheme $X$: part (2) gives surjectivity and part (1) injectivity of $M_g \to \mathcal{F}(D(g))$. So the map is an isomorphism on a basis.
The converse is clear. The equivalence follows with Proposition~\ref{schemes:proposition-tilde}(2). For exactness, kernels and cokernels
of maps $\widetilde M \to \widetilde N$ are the sheaves associated to the kernel and cokernel, since $\sim$ is exact. So a sequence of quasi-coherent sheaves
is isomorphic to $\sim$ applied to its sequence of global sections, and $\sim$ is exact and faithful: $\widetilde M = 0$ implies $M = 0$.
\end{proof}

\begin{corollary}\label{schemes:corollary-qc-properties}
On any scheme, kernels, cokernels, images, direct sums and tensor products of quasi-coherent sheaves are quasi-coherent, and pullbacks of
quasi-coherent sheaves are quasi-coherent. $\mathcal{F}$ is quasi-coherent if and only if $\mathcal{F}|_U \cong \mathcal{F}(U)^\sim$ for every affine open $U$.
On a Noetherian scheme the same statements hold for coherent sheaves (with pullbacks along morphisms of Noetherian schemes), and on
$\Spec A$ with $A$ Noetherian, $\widetilde M$ is coherent if and only if $M$ is finitely generated.
\end{corollary}

\begin{proof}
All statements are local, and for affine schemes they follow from Theorem~\ref{schemes:theorem-affine} and
Proposition~\ref{schemes:proposition-tilde}. For the characterisation of quasi-coherence on every affine open, restrict a cover with
$\mathcal{F}|_{U_i} = \widetilde{M_i}$ to the affine open $U$: then $\mathcal{F}|_U$ is quasi-coherent in the sense of the definition on $U$, and
Theorem~\ref{schemes:theorem-affine} applies. For coherence, submodules and quotients of finitely generated modules over a Noetherian ring are
finitely generated (Proposition~\ref{commutative-algebra:proposition-noetherian}). If $\widetilde M$ is coherent on $\Spec A$, then $M$ is locally, on a
finite cover by $D(f_i)$, generated by finitely many elements $m_{ij}/f_i^{n}$; the $m_{ij}$ generate $M$, since a module is zero when all its
localisations at the $f_i$ vanish.
\end{proof}

\begin{proposition}\label{schemes:proposition-pushforward}
Let $f \colon X \to Y$ be a morphism with $X$ Noetherian, or with $X$ quasi-compact and separated. If $\mathcal{F}$ is quasi-coherent on $X$, then $f_*\mathcal{F}$ is
quasi-coherent on $Y$.
\end{proposition}

\begin{proof}
The question is local on $Y$, so let $Y$ be affine. Cover $X$ by finitely many affine opens $U_i$. The intersections $U_{ij} = U_i \cap U_j$ are covered by
finitely many affine opens $U_{ijk}$ (in the separated case $U_{ij}$ itself is affine). By the sheaf axiom,
\[
0 \to f_*\mathcal{F} \to \bigoplus_if_*(\mathcal{F}|_{U_i}) \to \bigoplus_{i,j,k}f_*(\mathcal{F}|_{U_{ijk}})
\]
is exact. Each term on the right is the direct image of a quasi-coherent sheaf under a morphism of affine schemes, which is quasi-coherent by
Proposition~\ref{schemes:proposition-tilde}(3). So $f_*\mathcal{F}$ is a kernel of quasi-coherent sheaves.
\end{proof}

\begin{proposition}\label{schemes:proposition-closed-subschemes}
Let $X$ be a scheme. For a closed immersion $i \colon Y \to X$ the \emph{ideal sheaf} $\mathcal{I}_Y = \ker(\mathcal{O}_X \to i_*\mathcal{O}_Y)$ is quasi-coherent,
and $Y \mapsto \mathcal{I}_Y$ is a bijection from closed subschemes of $X$ to quasi-coherent ideal sheaves of $\mathcal{O}_X$. For $X = \Spec A$, every closed
subscheme is $\Spec A/I$ for a unique ideal $I$. If $X$ is Noetherian, $\mathcal{I}_Y$ is coherent.
\end{proposition}

\begin{proof}[Sketch of proof]
The question is local, so let $X = \Spec A$. The underlying space of $Y$ is closed in $X$, so $Y$ is quasi-compact. It is separated, since a closed immersion is
separated (Proposition~\ref{schemes:proposition-affine-separated}). By Proposition~\ref{schemes:proposition-pushforward}, $i_*\mathcal{O}_Y$ is
quasi-coherent, so the kernel $\mathcal{I}_Y = \widetilde I$ is quasi-coherent with $I = \Gamma(X, \mathcal{I}_Y)$. As $\mathcal{O}_X \to i_*\mathcal{O}_Y$ is onto,
$i_*\mathcal{O}_Y \cong (A/I)^\sim$, and its support $V(I)$ is the image of $Y$. Then $Y \to \Spec A/I$ is a homeomorphism with an isomorphism of structure
sheaves, so $Y \cong \Spec A/I$. Conversely, a quasi-coherent ideal $\widetilde I$ defines $\Spec A/I$, and on a general $X$ these glue. The omitted
verifications are the compatibilities of these constructions with restriction to smaller affine opens \cite[Corollary II.5.10]{hartshorne-ag}.
\end{proof}

\begin{definition}\label{schemes:definition-support}
The \emph{support} of a sheaf $\mathcal{F}$ is $\{x : \mathcal{F}_x \neq 0\}$. For a coherent sheaf it is closed: on $\Spec A$ with $M$ finitely generated,
$\mathrm{Supp}\,\widetilde M = V(\mathrm{Ann}\,M)$. The \emph{fibre} of $\mathcal{F}$ at $x$ is $\mathcal{F}(x) = \mathcal{F}_x \otimes k(x)$.
\end{definition}

\begin{proposition}\label{schemes:proposition-fibre-dimension}
Let $\mathcal{F}$ be coherent on a Noetherian scheme $X$.
\begin{enumerate}
\item (Nakayama) If sections $s_1, \ldots, s_r$ of $\mathcal{F}$ near $x$ span $\mathcal{F}(x)$, they generate $\mathcal{F}$ on a neighbourhood of $x$.
\item $x \mapsto \dim_{k(x)}\mathcal{F}(x)$ is upper semicontinuous.
\item If $X$ is reduced, $\mathcal{F}$ is locally free (locally isomorphic to $\mathcal{O}_X^r$) if and only if this function is locally constant.
\end{enumerate}
\end{proposition}

\begin{proof}
(1) By Nakayama's lemma (Corollary~\ref{commutative-algebra:corollary-nakayama-generators}) the $s_i$ generate the stalk $\mathcal{F}_x$. The cokernel of
$\mathcal{O}^r \to \mathcal{F}$ is coherent with $x$ outside its closed support. (2) follows from (1): generators at $x$ generate nearby, so the
fibre dimension cannot jump up near $x$. (3) If $X = \Spec A$ is reduced and the fibre dimension is $r$ near $x$, choose $s_1, \ldots, s_r$ as in (1) and
let $\mathcal{K}$ be the kernel of the resulting surjection $\mathcal{O}^r \to \mathcal{F}$. For each $y$ nearby, $\mathcal{O}^r(y) \to \mathcal{F}(y)$ is a
surjection of $r$-dimensional spaces, hence injective. So an element $(a_1, \ldots, a_r)$ of $\Gamma(\mathcal{K})$ has all $a_i$ vanishing at every point,
hence nilpotent, hence $0$. So $\mathcal{K} = 0$.
\end{proof}

\section{Locally free sheaves and the Picard group}\label{schemes:section-locally-free}

\begin{definition}\label{schemes:definition-locally-free}
An $\mathcal{O}_X$-module $\mathcal{E}$ is \emph{locally free of rank $r$} if it is locally isomorphic to $\mathcal{O}_X^r$, and \emph{invertible} if $r = 1$. Its
dual is $\mathcal{E}^\vee = \mathcal{H}om(\mathcal{E}, \mathcal{O}_X)$, its determinant $\det\mathcal{E} = \Lambda^r\mathcal{E}$. The \emph{Picard group} $\Pic(X)$ is
the group of isomorphism classes of invertible sheaves under $\otimes$, with inverse $\mathcal{L}^{-1} = \mathcal{L}^\vee$.
\end{definition}

\begin{proposition}\label{schemes:proposition-picard-cech}
$\Pic(X) \cong H^1(X, \mathcal{O}_X^*)$. Locally free sheaves of rank $r$ on $X$ correspond to isomorphism classes of \emph{algebraic vector bundles}: schemes
$E \to X$ locally isomorphic to $\Aa^r_U \to U$ with linear transition functions in $\mathrm{GL}_r(\mathcal{O}_X(U \cap V))$. The correspondence sends $E$ to its sheaf
of sections.
\end{proposition}

\begin{proof}
An invertible sheaf trivialised on an open cover $(U_i)$ by generators $e_i$ gives units $g_{ij} \in \mathcal{O}^*(U_i \cap U_j)$ with $e_j = g_{ij}e_i$. They form a
\v{C}ech cocycle, changing the $e_i$ changes it by a coboundary, and tensor product multiplies cocycles. Conversely a cocycle glues copies of
$\mathcal{O}_{U_i}$. So invertible sheaves trivialised on $(U_i)$ correspond to $\check H^1((U_i), \mathcal{O}^*)$, and passing to the limit over covers gives
$\check H^1(X, \mathcal{O}^*) = H^1(X, \mathcal{O}^*)$ (Remark~\ref{sheaves:remark-cech-limit}). The same cocycles with values in $\mathrm{GL}_r$ glue the
$\Aa^r_{U_i}$ to a vector bundle, as for manifolds (Proposition~\ref{vector-bundles:proposition-cocycles}).
\end{proof}

\section{Twisting sheaves on projective schemes}\label{schemes:section-twisting}

\begin{definition}\label{schemes:definition-twisting}
Let $S$ be a graded ring and $M$ a graded $S$-module. The sheaf $\widetilde M$ on $\Proj S$ has sections over $D_+(f)$ equal to $M_{(f)}$, the degree-zero part
of $M_f$, and stalks $M_{(\mathfrak{p})}$. For $n \in \ZZ$, $M(n)$ is $M$ with degrees shifted, $M(n)_d = M_{n+d}$, and $\mathcal{O}(n) = S(n)^\sim$. For an
$\mathcal{O}$-module $\mathcal{F}$ put $\mathcal{F}(n) = \mathcal{F} \otimes \mathcal{O}(n)$ and $\Gamma_*(\mathcal{F}) = \bigoplus_{n \in \ZZ}\Gamma(X, \mathcal{F}(n))$.
For a scheme $X$ with a closed immersion into $\PP^r_A$, $\mathcal{O}_X(1)$ is the restriction of $\mathcal{O}(1)$.
\end{definition}

\begin{proposition}\label{schemes:proposition-twisting}
Let $S$ be generated by $S_1$ as an $S_0$-algebra and $X = \Proj S$. Then $\mathcal{O}(n)$ is invertible, $\mathcal{O}(n) \otimes \mathcal{O}(m) \cong \mathcal{O}(n+m)$, and
$\widetilde M(n) \cong (M(n))^\sim$. For $S = A[x_0, \ldots, x_r]$ and $X = \PP^r_A$,
\[
\Gamma(X, \mathcal{O}(n)) \cong S_n \quad \text{for all } n \in \ZZ,
\]
the homogeneous polynomials of degree $n$ (zero for $n < 0$); so $\Gamma_*(\mathcal{O}_X) = S$.
\end{proposition}

\begin{proof}
Since $S$ is generated by $S_1$, the $D_+(f)$ with $f \in S_1$ cover $X$, and on $D_+(f)$ the module $S(n)_{(f)}$ is free over $S_{(f)}$ with basis $f^n$. This
proves invertibility, and the other isomorphisms follow similarly from $M(n)_{(f)} = M_{(f)} \cdot f^n$. For $\PP^r_A$, a global section of $\mathcal{O}(n)$ is a
compatible family of degree-$n$ elements $t_i \in S_{x_i}$. The $x_i$ are non-zero-divisors, so all $S_{x_i}$ embed in $S' = S_{x_0 \cdots x_r}$ and compatibility
means that the $t_i$ are one element $t$ of $\bigcap_iS_{x_i}$, whose homogeneous components lie in the intersection. An element of $S'$ is a sum of Laurent
monomials $x^a$, $a \in \ZZ^{r+1}$, with coefficients in $A$. It lies in $S_{x_i}$ exactly when only monomials with $a_j \geq 0$ for $j \neq i$ occur. Hence
$\bigcap_iS_{x_i} = S$, and the degree-$n$ part is $S_n$.
\end{proof}

\begin{theorem}[Serre]\label{schemes:theorem-serre-generation}
Let $X$ be a closed subscheme of $\PP^r_A$ with $A$ Noetherian, and $\mathcal{F}$ a coherent sheaf on $X$. Then there is $n_0$ such that for all $n \geq n_0$ the
sheaf $\mathcal{F}(n)$ is generated by finitely many global sections. Consequently $\mathcal{F}$ is a quotient of a finite direct sum of sheaves
$\mathcal{O}_X(-n)$.
\end{theorem}

\begin{proof}
Let $i \colon X \to \PP^r_A$ be the inclusion. Then $i_*\mathcal{F}$ is coherent (Proposition~\ref{schemes:proposition-pushforward}, with
finiteness checked on affines), and $i_*(\mathcal{F}(n)) = (i_*\mathcal{F})(n)$ since $\mathcal{O}_X(1) = i^*\mathcal{O}(1)$. So we may assume $X = \PP^r_A$. Cover
$X$ by $U_i = D_+(x_i) = X_{x_i}$, where $x_i \in \Gamma(X, \mathcal{O}(1))$, and write $\mathcal{F}|_{U_i} = \widetilde{M_i}$ with $M_i$ generated by finitely
many sections $s_{ij}$. By Lemma~\ref{schemes:lemma-extend-sections}(2) with $\mathcal{L} = \mathcal{O}(1)$ and $f = x_i$, for $n$ large each
$s_{ij} \otimes x_i^n$ extends to a global section of $\mathcal{F}(n)$, and we can take one $n$ for the finitely many sections. On $U_i$, $x_i^n$ generates
$\mathcal{O}(n)$, so the sections $s_{ij} \otimes x_i^n$ generate $\mathcal{F}(n)|_{U_i}$. The same holds for all larger $n$ (multiply by $x_i$). Finally, $N$ global
sections generating $\mathcal{F}(n)$ give a surjection $\mathcal{O}^N \to \mathcal{F}(n)$, that is, $\mathcal{O}(-n)^N \to \mathcal{F}$.
\end{proof}

\begin{definition}\label{schemes:definition-very-ample}
An invertible sheaf $\mathcal{L}$ on a scheme $X$ of finite type over $A$ is \emph{very ample} (relative to $A$) if $\mathcal{L} \cong i^*\mathcal{O}(1)$ for an
immersion $i \colon X \to \PP^r_A$, and \emph{ample} if for every coherent $\mathcal{F}$ the sheaf $\mathcal{F} \otimes \mathcal{L}^n$ is globally generated for
$n \gg 0$. By Theorem~\ref{schemes:theorem-serre-generation}, very ample sheaves on projective schemes over Noetherian rings are ample. The
converse relation, that some power of an ample sheaf is very ample, is in Chapter~\ref{divisors}.
\end{definition}

\section{Sheaves of differentials}\label{schemes:section-differentials}

\begin{definition}\label{schemes:definition-differentials}
Let $f \colon X \to Y$ be a morphism of schemes. The \emph{sheaf of relative differentials} $\Omega_{X/Y}$ is the quasi-coherent $\mathcal{O}_X$-module
obtained by gluing the sheaves $(\Omega_{B/A})^\sim$ on affine opens $\Spec B \subseteq f^{-1}(\Spec A)$, together with the derivation
$d \colon \mathcal{O}_X \to \Omega_{X/Y}$. The gluing is possible because $\Omega$ commutes with localisation
(Proposition~\ref{commutative-algebra:proposition-differentials}(5)). For $X$ smooth of relative dimension $n$ over $Y$, the \emph{canonical sheaf} is
$\omega_{X/Y} = \Lambda^n\Omega_{X/Y}$, and the \emph{tangent sheaf} is $\mathcal{T}_{X/Y} = \Omega_{X/Y}^\vee$.
\end{definition}

\begin{proposition}\label{schemes:proposition-differentials}
\begin{enumerate}
\item For morphisms $X \to Y \to Z$ there is an exact sequence $f^*\Omega_{Y/Z} \to \Omega_{X/Z} \to \Omega_{X/Y} \to 0$.
\item For a closed subscheme $Y \subseteq X$ over $Z$ with ideal sheaf $\mathcal{I}$ there is an exact sequence
$\mathcal{I}/\mathcal{I}^2 \to \Omega_{X/Z}|_Y \to \Omega_{Y/Z} \to 0$ (the \emph{conormal sequence}).
\item $\Omega$ commutes with base change: $\Omega_{X_{Y'}/Y'} \cong g'^*\Omega_{X/Y}$.
\item If $X \to Y$ is smooth of relative dimension $n$, then $\Omega_{X/Y}$ is locally free of rank $n$. In particular, a nonsingular variety of
dimension $n$ over an algebraically closed field has $\Omega_{X/k}$ locally free of rank $n$.
\end{enumerate}
\end{proposition}

\begin{proof}
(1), (2) are Proposition~\ref{commutative-algebra:proposition-differentials}(3), (4) on affine opens. (3) holds because
$\Omega_{(B \otimes_AA')/A'} \cong \Omega_{B/A} \otimes_AA'$, which follows from the universal property. (4) Locally
$B = A[x_1, \ldots, x_{n+r}]/(g_1, \ldots, g_r)$ with $\det(\partial g_i/\partial x_{n+j})$ invertible. By
Proposition~\ref{commutative-algebra:proposition-differentials}(2), (4), $\Omega_{B/A}$ is the quotient of the free module on $dx_1, \ldots, dx_{n+r}$ by the
relations $dg_i = \sum_j(\partial g_i/\partial x_j)dx_j = 0$. Since the $r \times r$ block for $dx_{n+1}, \ldots, dx_{n+r}$ is invertible, these relations express
$dx_{n+1}, \ldots, dx_{n+r}$ uniquely in terms of $dx_1, \ldots, dx_n$, and $\Omega_{B/A}$ is free on $dx_1, \ldots, dx_n$. The last statement uses
Proposition~\ref{schemes:proposition-smooth-nonsingular}.
\end{proof}

\begin{theorem}[Euler sequence]\label{schemes:theorem-euler-sequence}
On $X = \PP^r_A$ there is an exact sequence of locally free sheaves
\[
0 \to \Omega_{X/A} \to \mathcal{O}(-1)^{r+1} \to \mathcal{O}_X \to 0,
\]
where the second map sends the basis section $e_j$ to $x_j$. Consequently $\omega_{\PP^r_A/A} \cong \mathcal{O}(-r-1)$.
\end{theorem}

\begin{proof}
Write $\mathcal{O}(-1)^{r+1} = (S(-1)^{r+1})^\sim$ with basis $e_0, \ldots, e_r$ of degree $1$, and let $\pi$ be the map $e_j \mapsto x_j$. It is surjective: on $U_i$,
the section $e_i/x_i$ maps to $1$. On $U_i = \Spec A[x_0/x_i, \ldots, x_r/x_i]$, $\Omega$ is free on the $d(x_j/x_i)$, $j \neq i$. Define
\[
\varphi_i \colon \Omega|_{U_i} \to \mathcal{O}(-1)^{r+1}|_{U_i}, \qquad d(x_j/x_i) \mapsto \frac{1}{x_i^2}\left(x_ie_j - x_je_i\right) .
\]
The composite with $\pi$ is $0$. The kernel of $\pi$ on $U_i$ consists of the $\sum_ja_je_j$ ($a_j$ of degree $-1$) with $\sum_ja_jx_j = 0$, that is,
$a_i = -\sum_{j \neq i}a_jx_j/x_i$, and then $\sum_ja_je_j = \sum_{j \neq i}(a_jx_i)\,\varphi_i(d(x_j/x_i))$ with $a_jx_i$ of degree $0$. The images $\varphi_i(d(x_j/x_i))$
are independent, since the coefficient of $e_j$ is $1/x_i$. So $\varphi_i$ is an isomorphism onto $\ker\pi|_{U_i}$. On $U_i \cap U_k$, writing
$x_j/x_i = (x_j/x_k)(x_k/x_i)$ and $d(x_k/x_i) = -(x_k/x_i)^2d(x_i/x_k)$ shows $\varphi_k(d(x_j/x_i)) = \varphi_i(d(x_j/x_i))$ by a direct computation. So the $\varphi_i$
glue. Taking determinants in the exact sequence gives $\omega \otimes \mathcal{O} \cong \det(\mathcal{O}(-1)^{r+1}) = \mathcal{O}(-r-1)$.
\end{proof}

\begin{proposition}[Adjunction]\label{schemes:proposition-adjunction}
Let $Y \subseteq X$ be nonsingular varieties over an algebraically closed field, of dimensions $n - 1$ and $n$, with ideal sheaf $\mathcal{I}$. Then
$\mathcal{I}$ is invertible, $\mathcal{I}/\mathcal{I}^2 \cong \mathcal{I}|_Y$ is the conormal sheaf, the conormal sequence
$0 \to \mathcal{I}/\mathcal{I}^2 \to \Omega_X|_Y \to \Omega_Y \to 0$ is exact on the left, and
\[
\omega_Y \cong \omega_X \otimes \mathcal{I}^{-1}|_Y .
\]
For a nonsingular hypersurface $Y$ of degree $d$ in $\PP^n$, $\omega_Y \cong \mathcal{O}_Y(d - n - 1)$.
\end{proposition}

\begin{proof}
Near a point of $Y$, $\mathcal{O}_X$ is a regular local ring, hence a unique factorisation domain (Remark~\ref{commutative-algebra:remark-regular-domain}),
and the height-one prime of $Y$ is principal, generated by $g$. So $\mathcal{I}$ is invertible and $\mathcal{I}/\mathcal{I}^2 = \mathcal{I} \otimes \mathcal{O}_Y$. The map
$\mathcal{I}/\mathcal{I}^2 \to \Omega_X|_Y$ sends $g$ to $dg$. Both sides of the right-exact conormal sequence are locally free, of ranks $1$, $n$ and $n - 1$
(Proposition~\ref{schemes:proposition-differentials}(4)). A surjection of locally free sheaves of ranks $n$ and $n - 1$ has locally free
kernel of rank $1$, which is the image of $\mathcal{I}/\mathcal{I}^2$. A surjection between invertible sheaves is an isomorphism, since locally it is
multiplication by a function that generates the unit ideal. So the sequence is exact. Taking
determinants gives $\omega_X|_Y \cong \omega_Y \otimes \mathcal{I}/\mathcal{I}^2$. For a hypersurface of degree $d$, $\mathcal{I} = \mathcal{O}(-d)$ and $\omega_{\PP^n} = \mathcal{O}(-n-1)$.
\end{proof}

\section{Cohomology of Coherent Sheaves}\label{schemes:section-coherent-cohomology}
The sheaf cohomology of Section~\ref{sheaves:section-sheaf-cohomology} applied to quasi-coherent sheaves on schemes is computable. On affine schemes it
vanishes in positive degrees (Theorem~\ref{schemes:theorem-affine-vanishing}), so on separated schemes it is computed by \v{C}ech
complexes of affine covers. On projective space everything can be computed explicitly
(Theorem~\ref{schemes:theorem-projective-space}). For a projective scheme over a field the cohomology groups of coherent sheaves are
finite-dimensional, those of $\mathcal{F}(n)$ vanish for $n \gg 0$, and the Euler characteristic $\chi(\mathcal{F}(n))$ is a polynomial in $n$, the
Hilbert polynomial. Serre duality relates $H^i$ and $H^{n-i}$. For complex projective manifolds these groups coincide with the Dolbeault
cohomology of Section~\ref{complex-manifolds:section-dolbeault} by GAGA (Chapter~\ref{gaga}), and so compute the Hodge numbers. References:
\cite[Chapter III]{hartshorne-ag}, \cite{stacks-project}.

\noindent\textbf{Prerequisites.} Section~\ref{schemes:section-coherent-sheaves} (Quasi-coherent and Coherent Sheaves) and Section~\ref{sheaves:section-sheaf-cohomology}
(Sheaf Cohomology).

For an $A$-module $M$ and $f_1, \ldots, f_r \in A$, the \emph{\v{C}ech complex} $\check C^\bullet(f; M)$ is
$\prod_{i_0}M_{f_{i_0}} \to \prod_{i_0 < i_1}M_{f_{i_0}f_{i_1}} \to \cdots$ with the alternating differential of
Definition~\ref{sheaves:definition-cech}. It is the \v{C}ech complex of $\widetilde M$ for the cover of $\bigcup_iD(f_i)$ by the $D(f_i)$.

\section{Vanishing on affine schemes}\label{schemes:section-affine}

\begin{lemma}\label{schemes:lemma-cech-exact}
If $f_1, \ldots, f_r$ generate the unit ideal of $A$, then the augmented \v{C}ech complex $0 \to M \to \check C^\bullet(f; M)$ is exact for every $A$-module $M$.
\end{lemma}

\begin{proof}
A module that becomes zero after inverting each $f_j$ is zero, because every prime omits some $f_j$
(Proposition~\ref{commutative-algebra:proposition-local-global}). Localisation is exact, so it suffices to show that the complex becomes exact after
inverting $f_j$. After inverting $f_j$ it is the augmented \v{C}ech complex of $M_{f_j}$ for the elements $f_1, \ldots, f_r$ of $A_{f_j}$, one of which is a unit.
The operator $(k\alpha)_{i_0 \cdots i_{p-1}} = \alpha_{ji_0 \cdots i_{p-1}}$ is then defined, since $M_{f_jf_{i_0} \cdots} = M_{f_{i_0} \cdots}$, and
$dk + kd = \mathrm{id}$ as in the proof of Lemma~\ref{sheaves:lemma-cech-resolution}. So the localised complex is contractible.
\end{proof}

\begin{lemma}[Cartan]\label{schemes:lemma-cartan}
Let $X$ be a topological space and $\mathcal{B}$ a basis of quasi-compact open sets closed under finite intersections. Let $\mathcal{F}$ be a sheaf of
abelian groups such that $\check H^p(\mathcal{V}, \mathcal{F}) = 0$ for all $p > 0$, all $U \in \mathcal{B}$ and all finite covers $\mathcal{V}$ of $U$ by members of
$\mathcal{B}$. Then $H^p(U, \mathcal{F}) = 0$ for all $U \in \mathcal{B}$ and $p > 0$.
\end{lemma}

\begin{proof}
Embed $\mathcal{F}$ into an injective sheaf $\mathcal{I}$ (Theorem~\ref{sheaves:theorem-enough-injectives}) and let $\mathcal{Q} = \mathcal{I}/\mathcal{F}$.

\emph{Step 1: $\mathcal{I}(U) \to \mathcal{Q}(U)$ is onto for $U \in \mathcal{B}$.} A section $s \in \mathcal{Q}(U)$ lifts locally, and by quasi-compactness we find a
finite cover of $U$ by $U_1, \ldots, U_m \in \mathcal{B}$ and lifts $t_i \in \mathcal{I}(U_i)$. Then $t_i - t_j \in \mathcal{F}(U_i \cap U_j)$ is a \v{C}ech $1$-cocycle, so by
hypothesis $t_i - t_j = c_j - c_i$ with $c_i \in \mathcal{F}(U_i)$. The sections $t_i + c_i$ agree on overlaps and glue to a lift of $s$.

\emph{Step 2: $\mathcal{Q}$ satisfies the hypothesis.} For a finite cover $\mathcal{V}$ of $U \in \mathcal{B}$ by members of $\mathcal{B}$, all intersections lie in
$\mathcal{B}$, so by Step 1 the sequence $0 \to \check C^\bullet(\mathcal{V}, \mathcal{F}) \to \check C^\bullet(\mathcal{V}, \mathcal{I}) \to \check C^\bullet(\mathcal{V}, \mathcal{Q}) \to 0$
is exact. Injective sheaves are flasque (Lemma~\ref{sheaves:lemma-injective-flasque}), so $\check H^p(\mathcal{V}, \mathcal{I}) = 0$ for $p > 0$
(Lemma~\ref{sheaves:lemma-cech-resolution}). The long exact sequence gives $\check H^p(\mathcal{V}, \mathcal{Q}) \cong \check H^{p+1}(\mathcal{V}, \mathcal{F}) = 0$.

\emph{Step 3.} By the long exact sequence of $0 \to \mathcal{F} \to \mathcal{I} \to \mathcal{Q} \to 0$ on $U$, $H^1(U, \mathcal{F}) = \operatorname{coker}(\mathcal{I}(U) \to \mathcal{Q}(U)) = 0$
by Step 1, and $H^{p+1}(U, \mathcal{F}) \cong H^p(U, \mathcal{Q})$ for $p \geq 1$. By Step 2 and induction on $p$, applied to all sheaves satisfying the hypothesis at
once, $H^p(U, \mathcal{F}) = 0$ for all $p \geq 1$.
\end{proof}

\begin{theorem}[Serre]\label{schemes:theorem-affine-vanishing}
Let $X = \Spec A$ and $\mathcal{F}$ a quasi-coherent sheaf on $X$. Then $H^p(X, \mathcal{F}) = 0$ for all $p > 0$.
\end{theorem}

\begin{proof}
Take $\mathcal{B}$ to be the basic open sets $D(g)$. They are quasi-compact (Proposition~\ref{schemes:proposition-spec-topology}) and
$D(g) \cap D(h) = D(gh)$. A finite cover of $D(g) \cong \Spec A_g$ by basic open sets $D(f_i)$ with $D(f_i) \subseteq D(g)$ is a cover of $\Spec A_g$ by the $D(f_i)$, and the
$f_i$ generate the unit ideal of $A_g$. By Theorem~\ref{schemes:theorem-affine}, $\mathcal{F} = \widetilde M$ and the \v{C}ech complex of the cover is
$\check C^\bullet(f; M_g)$, which is exact in positive degrees by Lemma~\ref{schemes:lemma-cech-exact}. Apply
Lemma~\ref{schemes:lemma-cartan}.
\end{proof}

\begin{counterexample}\label{schemes:counterexample-not-quasi-coherent}
Theorem~\ref{schemes:theorem-affine-vanishing} needs quasi-coherence. Let $X = \Aa^1_k = \Spec k[t]$ with its Zariski topology, and let $0$, $1$ be the corresponding closed points. Consider the morphism of sheaves of
abelian groups
\[
\ZZ_X \to (i_0)_*\ZZ \oplus (i_1)_*\ZZ
\]
to the skyscraper sheaves at $0$ and $1$, taking the values of a locally constant function at the two points, and let $\mathcal{K}$ be its kernel. It is surjective on stalks: the stalk of $(i_0)_*\ZZ$ at a point $y$ is $\ZZ$
if every open neighbourhood of $y$ contains $0$, that is, if $y = 0$, since $\{0\}$ is closed, and $0$ otherwise; and at $y = 0$ the map on stalks is the identity of $\ZZ$. So
$0 \to \mathcal{K} \to \ZZ_X \to (i_0)_*\ZZ \oplus (i_1)_*\ZZ \to 0$ is exact. Since $X$ is irreducible, every nonempty open subset is irreducible, hence connected, so $\ZZ_X(U) = \ZZ$ for all nonempty $U$, with the identity as
restriction maps; thus $\ZZ_X$ is flasque and $H^1(X, \ZZ_X) = 0$ (Proposition~\ref{sheaves:proposition-flasque}), and $\Gamma(X, \ZZ_X) = \ZZ$ maps to $\ZZ \oplus \ZZ$ diagonally. The long exact sequence gives
\[
H^1(X, \mathcal{K}) \cong (\ZZ \oplus \ZZ)/\ZZ \cong \ZZ \neq 0,
\]
although $X$ is affine. The sheaf $\mathcal{K}$, the extension by zero of $\ZZ$ from $X \setminus \{0, 1\}$ (Definition~\ref{sheaves:definition-extension-by-zero}), is not a sheaf of $\mathcal{O}_X$-modules.
\end{counterexample}

\begin{corollary}\label{schemes:corollary-cech-computes}
Let $X$ be a quasi-compact separated scheme, $\mathcal{U}$ a finite cover by affine open subsets and $\mathcal{F}$ quasi-coherent. Then
$H^p(X, \mathcal{F}) \cong \check H^p(\mathcal{U}, \mathcal{F})$. In particular $H^p(X, \mathcal{F}) = 0$ for $p \geq |\mathcal{U}|$.
\end{corollary}

\begin{proof}
Finite intersections of affine opens in a separated scheme are affine (as in the proof of Lemma~\ref{schemes:lemma-extend-sections}), so
Theorem~\ref{schemes:theorem-affine-vanishing} and Leray's theorem (Theorem~\ref{sheaves:theorem-leray}) apply. The alternating \v{C}ech
complex of a cover with $m$ members vanishes in degrees $\geq m$.
\end{proof}

\begin{remark}\label{schemes:remark-serre-criterion}
Conversely, Serre's criterion says that a Noetherian scheme on which $H^1$ of every quasi-coherent (even every coherent) ideal sheaf vanishes is
affine \cite[Theorem III.3.7]{hartshorne-ag}. Grothendieck's vanishing theorem says that $H^p(X, \mathcal{F}) = 0$ for $p > \dim X$ for every sheaf of abelian
groups on a Noetherian space $X$ \cite[Theorem III.2.7]{hartshorne-ag}.
\end{remark}

\section{The cohomology of projective space}\label{schemes:section-projective-space}

\begin{theorem}\label{schemes:theorem-projective-space}
Let $A$ be a Noetherian ring, $S = A[x_0, \ldots, x_r]$ and $X = \PP^r_A$ with $r \geq 1$.
\begin{enumerate}
\item $\bigoplus_nH^0(X, \mathcal{O}(n)) \cong S$.
\item $H^i(X, \mathcal{O}(n)) = 0$ for $0 < i < r$ and all $n$.
\item $H^r(X, \mathcal{O}(n))$ is the free $A$-module with basis the Laurent monomials $x_0^{a_0} \cdots x_r^{a_r}$ with all $a_i \leq -1$ and $\sum a_i = n$. In particular
$H^r(X, \mathcal{O}(-r-1)) \cong A$, spanned by $(x_0 \cdots x_r)^{-1}$, and $H^r(X, \mathcal{O}(n)) = 0$ for $n > -r-1$.
\item The multiplication pairing $H^0(X, \mathcal{O}(n)) \times H^r(X, \mathcal{O}(-n-r-1)) \to H^r(X, \mathcal{O}(-r-1)) \cong A$ is perfect.
\end{enumerate}
\end{theorem}

\begin{proof}
Let $\mathcal{F} = \bigoplus_n\mathcal{O}(n)$, a quasi-coherent sheaf. By Corollary~\ref{schemes:corollary-cech-computes} with the standard cover
$U_i = D_+(x_i)$, $H^i(X, \mathcal{O}(n))$ is the degree-$n$ part of the cohomology of the graded complex
\[
\prod_{i_0}S_{x_{i_0}} \to \prod_{i_0 < i_1}S_{x_{i_0}x_{i_1}} \to \cdots \to S_{x_0 \cdots x_r},
\]
which is the \v{C}ech complex of $\mathcal{F}$. (1) is Proposition~\ref{schemes:proposition-twisting}. (3) $H^r$ is the cokernel of $\prod_kS_{x_0 \cdots \widehat{x_k} \cdots x_r} \to S_{x_0 \cdots x_r}$. The
target is free on all Laurent monomials, and the image is spanned by those in which some exponent $a_k$ is $\geq 0$. The monomials with all $a_i \leq -1$ give a
basis of the cokernel. (4) In the basis of (3), $x^b \cdot x^{a}$ with $b_i \geq 0$, $\sum b_i = n$, and $a_i \leq -1$, $\sum a_i = -n-r-1$ is zero in $H^r$ unless
$b_i + a_i = -1$ for all $i$. So the monomial bases $\{x^b\}$ and $\{x^{-b-\mathbf{1}}\}$ are dual.

(2) By induction on $r$. The case $r = 1$ is empty. Let $r \geq 2$ and $M^i = \bigoplus_nH^i(X, \mathcal{O}(n))$. \emph{Every element of $M^i$, $i > 0$, is killed by a power of
$x_r$.} Indeed, localising the \v{C}ech complex at $x_r$ gives the \v{C}ech complex of $S_{x_r}$ for the cover of $U_r$ by the affine sets $U_i \cap U_r$. Its
cohomology is $H^i(U_r, \mathcal{F}|_{U_r}) = 0$ for $i > 0$, as $U_r$ is affine (Theorem~\ref{schemes:theorem-affine-vanishing}), and localisation is exact.
\emph{Multiplication by $x_r$ is injective on $M^i$ for $0 < i < r$.} Let $H = V_+(x_r) \cong \PP^{r-1}_A$. The exact sequence
$0 \to S(-1) \xrightarrow{x_r} S \to S/(x_r) \to 0$ gives $0 \to \mathcal{O}(n-1) \to \mathcal{O}(n) \to \mathcal{O}_H(n) \to 0$, with $H^i(X, \mathcal{O}_H(n)) = H^i(H, \mathcal{O}_H(n))$
(Example~\ref{sheaves:example-leray}). In the long exact sequence
\[
H^{i-1}(H, \mathcal{O}_H(n)) \to H^i(X, \mathcal{O}(n-1)) \xrightarrow{x_r} H^i(X, \mathcal{O}(n)),
\]
the left term vanishes for $2 \leq i \leq r - 1$ by induction ($0 < i - 1 < r - 1$). For $i = 1$, the map $H^0(X, \mathcal{O}(n)) = S_n \to H^0(H, \mathcal{O}_H(n)) = (S/x_r)_n$ is
onto, so again $x_r$ is injective on $H^1$. An element killed by a power of $x_r$, on which $x_r$ is injective, is zero. So $M^i = 0$ for $0 < i < r$.
\end{proof}

\begin{corollary}\label{schemes:corollary-p1}
On $\PP^1_k$, $h^0(\mathcal{O}(n)) = \max(n + 1, 0)$ and $h^1(\mathcal{O}(n)) = \max(-n - 1, 0)$. So $\chi(\mathcal{O}(n)) = n + 1$ for all $n$, and $H^1(\PP^1, \mathcal{O}(-2)) \cong k$.
\end{corollary}

\begin{proof}
Count the monomials $x_0^{a_0}x_1^{a_1}$ with $a_0 + a_1 = n$ and $a_i \geq 0$, respectively $a_i \leq -1$.
\end{proof}

\section{Finiteness and Serre vanishing}\label{schemes:section-finiteness}

\begin{theorem}[Serre]\label{schemes:theorem-finiteness}
Let $X$ be a projective scheme over a Noetherian ring $A$ with a very ample sheaf $\mathcal{O}_X(1)$, and $\mathcal{F}$ a coherent sheaf on $X$. Then:
\begin{enumerate}
\item $H^i(X, \mathcal{F})$ is a finitely generated $A$-module for every $i$;
\item there is $n_0$ with $H^i(X, \mathcal{F}(n)) = 0$ for all $i > 0$ and $n \geq n_0$.
\end{enumerate}
\end{theorem}

\begin{proof}
Pushing forward along the closed immersion $X \to \PP^r_A$ changes neither the cohomology nor coherence (Example~\ref{sheaves:example-leray}), so
let $X = \PP^r_A$. Both statements hold for $i > r$ by Corollary~\ref{schemes:corollary-cech-computes}. We use descending induction on $i$, for all
coherent sheaves simultaneously. By Theorem~\ref{schemes:theorem-serre-generation} there is an exact sequence $0 \to \mathcal{R} \to \mathcal{E} \to \mathcal{F} \to 0$
with $\mathcal{E}$ a finite direct sum of sheaves $\mathcal{O}(q_j)$ and $\mathcal{R}$ coherent (Corollary~\ref{schemes:corollary-qc-properties}). By
Theorem~\ref{schemes:theorem-projective-space}, $H^i(\mathcal{E})$ is finitely generated and $H^i(\mathcal{E}(n)) = 0$ for $i > 0$ and $n$ large. The exact
sequence $H^i(\mathcal{E}(n)) \to H^i(\mathcal{F}(n)) \to H^{i+1}(\mathcal{R}(n))$ and the induction hypothesis for $\mathcal{R}$ give both statements for $\mathcal{F}$ in degree
$i$. For (1), $A$ is Noetherian, so an extension of a submodule of a finitely generated module by a quotient of a finitely generated module is
finitely generated.
\end{proof}

\begin{corollary}\label{schemes:corollary-finite-dimensional}
For a projective scheme $X$ over a field $k$ and $\mathcal{F}$ coherent, all $H^i(X, \mathcal{F})$ are finite-dimensional, and they vanish for $i > \dim X$. In particular
$\Gamma(X, \mathcal{O}_X)$ is a finite-dimensional $k$-algebra, equal to $k$ if $X$ is geometrically integral.
\end{corollary}

\begin{proof}
Finiteness is Theorem~\ref{schemes:theorem-finiteness}(1), and vanishing above the dimension is Grothendieck's theorem
(Remark~\ref{schemes:remark-serre-criterion}). If $X$ is integral, $\Gamma(X, \mathcal{O}_X)$ is a domain finite over $k$, hence a field finite over $k$. If $X$
is geometrically integral, its base change to $\bar k$ is $\Gamma(X_{\bar k}, \mathcal{O}) = \bar k$ by Theorem~\ref{varieties:theorem-projective-constant}
(cohomology commutes with flat base change of fields, by the \v{C}ech description), so $\Gamma(X, \mathcal{O}_X) = k$.
\end{proof}

\section{Euler characteristic and Hilbert polynomial}\label{schemes:section-hilbert}

\begin{definition}\label{schemes:definition-euler}
For a coherent sheaf $\mathcal{F}$ on a projective scheme $X$ over a field $k$, the \emph{Euler characteristic} is $\chi(\mathcal{F}) = \sum_i(-1)^ih^i(\mathcal{F})$, where
$h^i = \dim_kH^i$. It is additive in short exact sequences, by the long exact sequence.
\end{definition}

\begin{theorem}\label{schemes:theorem-hilbert-polynomial}
Let $X \subseteq \PP^r_k$ be a projective scheme over an infinite field $k$ and $\mathcal{F}$ a coherent sheaf on $X$. There is a polynomial $P_\mathcal{F} \in \QQ[t]$, the
\emph{Hilbert polynomial}, with $\chi(\mathcal{F}(n)) = P_\mathcal{F}(n)$ for all $n \in \ZZ$. Its degree is $\dim\mathrm{Supp}\,\mathcal{F}$ ($-\infty$ if $\mathcal{F} = 0$), and
$h^0(\mathcal{F}(n)) = P_\mathcal{F}(n)$ for $n \gg 0$.
\end{theorem}

\begin{proof}
Induction on $d = \dim\mathrm{Supp}\,\mathcal{F}$. If $d \leq 0$, the support is finite, $\mathcal{O}(1)$ is trivial on a neighbourhood of it, so $\mathcal{F}(n) \cong \mathcal{F}$ and
$\chi$ is constant, equal to $h^0(\mathcal{F})$, which is nonzero if $\mathcal{F} \neq 0$ (a sheaf with finite support is a sum of skyscrapers). Let $d \geq 1$. The
\emph{associated points} of $\mathcal{F}$, the points $x$ such that $\mathfrak{m}_x$ is an associated prime of $\mathcal{F}_x$, are finite in number
(Proposition~\ref{commutative-algebra:proposition-associated-primes}, on a finite affine cover). Since $k$ is infinite, there is a linear form $h$ vanishing at none of
them: $S_1$ is not a finite union of proper subspaces. Then multiplication by $h$ is injective on $\mathcal{F}(-1) \to \mathcal{F}$, because its kernel would have
an associated point among those of $\mathcal{F}$. So there is an exact sequence
\[
0 \to \mathcal{F}(n-1) \xrightarrow{h} \mathcal{F}(n) \to \mathcal{F}_H(n) \to 0, \qquad \mathcal{F}_H = \mathcal{F}/h\mathcal{F}.
\]
$\mathrm{Supp}\,\mathcal{F}_H = \mathrm{Supp}\,\mathcal{F} \cap V_+(h)$ has dimension $d - 1$ by Theorem~\ref{varieties:theorem-projective-dimension}, since $h$ vanishes at no generic
point of a component of the support (these are associated points). By induction $\chi(\mathcal{F}(n)) - \chi(\mathcal{F}(n-1)) = P_{\mathcal{F}_H}(n)$ is a polynomial of degree
$d - 1$, so $\chi(\mathcal{F}(n))$ is a polynomial of degree $d$. The last statement follows from Theorem~\ref{schemes:theorem-finiteness}(2).
\end{proof}

\begin{definition}\label{schemes:definition-degree-genus}
For a closed subscheme $X \subseteq \PP^r_k$ of dimension $d$ with Hilbert polynomial $P_X = P_{\mathcal{O}_X}$, the \emph{degree} of $X$ is $d!$ times the leading
coefficient of $P_X$, and the \emph{arithmetic genus} is $p_a(X) = (-1)^d(P_X(0) - 1)$.
\end{definition}

\begin{example}\label{schemes:example-hilbert-polynomials}
\begin{enumerate}
\item $P_{\PP^r}(n) = \binom{n+r}{r}$ for all $n$, by Theorem~\ref{schemes:theorem-projective-space}. So $\PP^r$ has degree $1$ and $p_a = 0$.
\item For a hypersurface $X = V_+(f)$ of degree $d$, the sequence $0 \to \mathcal{O}(n-d) \to \mathcal{O}(n) \to \mathcal{O}_X(n) \to 0$ gives
$P_X(n) = \binom{n+r}{r} - \binom{n-d+r}{r}$, of degree $r - 1$ with leading coefficient $d/(r-1)!$. So $X$ has degree $d$.
\item For a plane curve of degree $d$, $P_X(n) = dn + 1 - \frac{(d-1)(d-2)}{2}$, so $p_a = \frac{(d-1)(d-2)}{2}$: $0$ for lines and conics, $1$ for cubics, $3$ for quartics.
\item The twisted cubic $C = v_3(\PP^1) \subseteq \PP^3$ has $\mathcal{O}_C(n) = \mathcal{O}_{\PP^1}(3n)$, so $P_C(n) = 3n + 1$: degree $3$, genus $0$.
\end{enumerate}
\end{example}

\section{Serre duality}\label{schemes:section-serre-duality}

\begin{theorem}[Serre duality on projective space]\label{schemes:theorem-duality-projective}
Let $X = \PP^r_k$ and $\omega = \mathcal{O}(-r-1)$, with the isomorphism $H^r(X, \omega) \cong k$ of Theorem~\ref{schemes:theorem-projective-space}. For every coherent
$\mathcal{F}$ the pairing $\Hom(\mathcal{F}, \omega) \times H^r(X, \mathcal{F}) \to H^r(X, \omega) \cong k$ is perfect.
\end{theorem}

\begin{proof}
Both $\mathcal{F} \mapsto \Hom(\mathcal{F}, \omega)$ and $\mathcal{F} \mapsto H^r(X, \mathcal{F})^\vee$ are contravariant and left exact: the second because $H^{r+1} = 0$, so $H^r$ is right exact.
The pairing gives a natural transformation between them, which is an isomorphism for $\mathcal{F} = \mathcal{O}(q)$ by
Theorem~\ref{schemes:theorem-projective-space}(4), since $\Hom(\mathcal{O}(q), \omega) = H^0(\mathcal{O}(-q-r-1))$, and hence for finite sums of such.
Every coherent $\mathcal{F}$ has a presentation $\mathcal{E}_1 \to \mathcal{E}_0 \to \mathcal{F} \to 0$ by finite sums of twisting sheaves
(Theorem~\ref{schemes:theorem-serre-generation}, applied twice). By the five lemma the transformation is an isomorphism for $\mathcal{F}$.
\end{proof}

\begin{theorem}[Serre duality]\label{schemes:theorem-serre-duality}
Let $X$ be a nonsingular projective variety of dimension $n$ over an algebraically closed field $k$, with canonical sheaf $\omega_X = \Lambda^n\Omega_{X/k}$. There is an
isomorphism $H^n(X, \omega_X) \cong k$, and for every locally free sheaf $\mathcal{E}$ and every $i$ the pairing
\[
H^i(X, \mathcal{E}) \times H^{n-i}(X, \mathcal{E}^\vee \otimes \omega_X) \to H^n(X, \omega_X) \cong k
\]
is perfect. More generally $\mathrm{Ext}^i(\mathcal{F}, \omega_X) \cong H^{n-i}(X, \mathcal{F})^\vee$ for coherent $\mathcal{F}$.
\end{theorem}

This is \cite[Theorem III.7.6 and Corollary III.7.7]{hartshorne-ag}, where it is deduced from Theorem~\ref{schemes:theorem-duality-projective} and a computation of
$\mathrm{Ext}$ sheaves for the embedding $X \subseteq \PP^r$, and \cite[Corollary III.7.12]{hartshorne-ag} for the identification of the dualising sheaf with $\omega_X$. For
$k = \CC$, via GAGA (Chapter~\ref{gaga}), it agrees with the analytic Serre duality of Theorem~\ref{complex-manifolds:theorem-serre-duality}.

\begin{corollary}\label{schemes:corollary-hodge-symmetry-serre}
For a nonsingular projective variety over $k$ put $h^{p,q} = \dim H^q(X, \Omega^p_X)$. Then $h^{p,q} = h^{n-p,n-q}$, and $\chi(\omega_X) = (-1)^n\chi(\mathcal{O}_X)$.
\end{corollary}

\begin{proof}
$\Omega^p$ is locally free, and $(\Omega^p)^\vee \otimes \omega \cong \Omega^{n-p}$ by the perfect pairing $\Omega^p \otimes \Omega^{n-p} \to \omega$. Apply
Theorem~\ref{schemes:theorem-serre-duality}. For the second, take $\mathcal{E} = \mathcal{O}_X$.
\end{proof}

\section{Semicontinuity and base change}\label{schemes:section-base-change-coherent-cohomology}

\begin{theorem}[Semicontinuity]\label{schemes:theorem-semicontinuity}
Let $f \colon X \to Y$ be a projective morphism of Noetherian schemes and $\mathcal{F}$ a coherent sheaf on $X$, flat over $Y$. For $y \in Y$ let
$\mathcal{F}_y$ be the restriction to the fibre $X_y$. Then:
\begin{enumerate}
\item $y \mapsto h^i(X_y, \mathcal{F}_y)$ is upper semicontinuous for every $i$;
\item $y \mapsto \chi(\mathcal{F}_y)$ is locally constant, and so is the Hilbert polynomial of $\mathcal{F}_y$;
\item (cohomology and base change) if $Y$ is integral and $h^i(X_y, \mathcal{F}_y)$ is constant, then $R^if_*\mathcal{F}$ is locally free and
$R^if_*\mathcal{F} \otimes k(y) \cong H^i(X_y, \mathcal{F}_y)$ for all $y$.
\end{enumerate}
\end{theorem}

This is \cite[Theorem III.12.8, Corollary III.12.9 and Theorem III.9.9]{hartshorne-ag}. The proof represents the cohomology of all fibres by one bounded complex
of finitely generated locally free sheaves on $Y$. The theorem implies that Hodge numbers are upper semicontinuous in algebraic families; in the
K\"ahler case they are in fact constant, by the Hodge decomposition and the constancy of Betti numbers in smooth families (Chapter~\ref{vhs}).

\begin{counterexample}\label{schemes:counterexample-jump}
Individual $h^i$ can jump, although $\chi$ cannot. For the family of line bundles $\mathcal{L}_t$ of degree $0$ on a fixed elliptic curve $E$, parametrised
by $t \in E$ (the Poincar\'e bundle on $E \times E$), one has $h^0(\mathcal{L}_t) = 1$ if $\mathcal{L}_t$ is trivial and $0$ otherwise, while
$\chi(\mathcal{L}_t) = 0$ is constant (Chapter~\ref{curves}).
\end{counterexample}

\section{Exercises}\label{schemes:section-exercises}

\subsection{Schemes}\label{schemes:subsection-exercises-schemes}

\begin{exercise}\label{schemes:exercise-spec-z-x}
Describe the points of $\Spec\ZZ[x]$: show that they are $(0)$, the primes $(f)$ with $f$ irreducible, and the maximal ideals $(p, f)$ with $p$ prime and
$f$ irreducible modulo $p$.
\end{exercise}

\begin{solution}
Let $\mathfrak{p}$ be a prime. If $\mathfrak{p} \cap \ZZ = (p)$ with $p$ prime, then $\mathfrak{p}$ corresponds to a prime of $\mathbb{F}_p[x]$, a principal ideal domain,
namely $0$ or $(\bar f)$ with $\bar f$ irreducible. This gives $(p)$ and $(p, f)$, and the latter is maximal. If $\mathfrak{p} \cap \ZZ = 0$, then $\mathfrak{p}$
corresponds to a prime of $\QQ[x]$ (localisation at $\ZZ \setminus 0$), which is $0$ or $(g)$ with $g$ irreducible. Then $\mathfrak{p} = g\QQ[x] \cap \ZZ[x] = (f)$,
where $f$ is $g$ made primitive (Gauss's lemma, Lemma~\ref{rings:lemma-gauss}). The primes $(p)$ are the primes $(f)$ with $f = p$ constant.
\end{solution}

\begin{exercise}\label{schemes:exercise-dual-numbers}
Let $X$ be a scheme over an algebraically closed field $k$ and $D = \Spec k[\varepsilon]/(\varepsilon^2)$. Show that $X(D)$ is in bijection with pairs of a closed
point $x \in X(k)$ and a vector of $(\mathfrak{m}_x/\mathfrak{m}_x^2)^\vee$.
\end{exercise}

\begin{solution}
The image of $D$ is a point $x$, and the morphism is a local $k$-algebra map $\varphi \colon \mathcal{O}_{X,x} \to k[\varepsilon]$. Composing with
$k[\varepsilon] \to k$ gives $k(x) = k$. Write $\varphi(a) = a(x) + \delta(a)\varepsilon$. Then $\delta$ is a $k$-linear derivation at $x$ that kills $k$ and $\mathfrak{m}_x^2$,
since $\varepsilon^2 = 0$. So $\delta$ is a linear form on $\mathfrak{m}_x/\mathfrak{m}_x^2$, and conversely every such form defines a local homomorphism.
\end{solution}

\begin{exercise}\label{schemes:exercise-dvr}
Let $R$ be a discrete valuation ring with fraction field $K$ and residue field $\kappa$. Describe $\Spec R$ and the fibres of $X \to \Spec R$ for
$X = \Spec R[x, y]/(xy - \pi)$, where $\pi$ is a uniformiser.
\end{exercise}

\begin{solution}
$\Spec R$ has two points: the closed point $\mathfrak{m}$ and the generic point $(0)$, which is open. The generic fibre is the smooth conic $xy = \pi$ over
$K$, isomorphic to $\Aa^1_K \setminus \{0\}$. The special fibre is $xy = 0$ over $\kappa$, two lines meeting in a point. $X$ is a domain dominating $\Spec R$,
so the morphism is flat (Proposition~\ref{schemes:proposition-flat}).
\end{solution}

\begin{exercise}\label{schemes:exercise-proj-conic}
Show that $\Proj k[x, y, z]/(xz - y^2) \cong \PP^1_k$ for every field $k$.
\end{exercise}

\begin{solution}
The ring $S = k[x, y, z]/(xz - y^2)$ is isomorphic to the subring $k[s^2, st, t^2] \subseteq k[s, t]$ of polynomials of even degree, with degrees halved.
For a graded ring $T$ and $T^{(2)} = \bigoplus_dT_{2d}$ one has $\Proj T^{(2)} \cong \Proj T$, since $D_+(f^2)$ and $D_+(f)$ have the same ring
$T_{(f)} = T^{(2)}_{(f^2)}$ \cite[Exercise II.5.13]{hartshorne-ag}. Hence $\Proj S \cong \Proj k[s, t] = \PP^1_k$.
\end{solution}

\subsection{Quasi-coherent and Coherent Sheaves}\label{schemes:subsection-exercises-coherent-sheaves}

\begin{exercise}\label{schemes:exercise-not-quasi-coherent}
Let $X = \Spec R$ for a discrete valuation ring $R$ with fraction field $K$, and let $\mathcal{F}$ be the sheaf with $\mathcal{F}(X) = 0$ and $\mathcal{F}(\{\eta\}) = K$ on
the open generic point. Show that $\mathcal{F}$ is an $\mathcal{O}_X$-module that is not quasi-coherent, while $j_*K$ for $j \colon \{\eta\} \to X$ is quasi-coherent.
\end{exercise}

\begin{solution}
If $\mathcal{F}$ were quasi-coherent, it would be $\widetilde{\Gamma(X, \mathcal{F})} = 0$ (Theorem~\ref{schemes:theorem-affine}), but $\mathcal{F}_\eta = K \neq 0$.
On the other hand $j_*K = \widetilde K$, since $\widetilde K(\{\eta\}) = K_{(0)} = K$ and $\widetilde K(X) = K$.
\end{solution}

\begin{exercise}\label{schemes:exercise-o-minus-one}
Show that $\mathcal{O}(-1)$ on $\PP^r_k$ is the sheaf of sections of the tautological line bundle
$\{(x, \ell) \in \Aa^{r+1} \times \PP^r : x \in \ell\}$, and that $\mathcal{O}(1)$ has no nonzero section vanishing on a hyperplane twice.
\end{exercise}

\begin{solution}
On $U_i$ the tautological bundle is trivialised by the section $\ell \mapsto (x_0/x_i, \ldots, x_r/x_i)$, and on $U_i \cap U_j$ these trivialisations differ by
$x_i/x_j$. This is the cocycle of $\mathcal{O}(-1)$, whose generator on $U_i$ is $1/x_i$: $1/x_j = (x_i/x_j)(1/x_i)$. A section of $\mathcal{O}(1)$ is a linear form
by Proposition~\ref{schemes:proposition-twisting}, and a nonzero linear form vanishes on its hyperplane to order exactly $1$.
\end{solution}

\begin{exercise}\label{schemes:exercise-canonical-curves}
Using Proposition~\ref{schemes:proposition-adjunction}, show that a nonsingular plane cubic has trivial canonical sheaf and that a nonsingular plane
quartic has $\omega \cong \mathcal{O}(1)$, so that $\Gamma(Y, \omega_Y)$ contains the three linear forms. (In fact it is spanned by them, and the genus is
$3$; see Chapter~\ref{curves}.)
\end{exercise}

\begin{solution}
$d - n - 1 = 3 - 3 = 0$ and $4 - 3 = 1$. The restriction $\Gamma(\PP^2, \mathcal{O}(1)) \to \Gamma(Y, \mathcal{O}_Y(1))$ is injective because a linear form vanishing on $Y$
would be divisible by the quartic equation. So $\Gamma(Y, \omega_Y)$ contains a $3$-dimensional space.
\end{solution}

\begin{exercise}\label{schemes:exercise-differentials-cusp}
Compute $\Omega_{X/k}$ for the cusp $X = \Spec k[x, y]/(y^2 - x^3)$ and show that it is not locally free.
\end{exercise}

\begin{solution}
By Example~\ref{commutative-algebra:example-cusp}, $\Omega = (B\,dx \oplus B\,dy)/(2y\,dy - 3x^2dx)$. Its fibre at the origin has dimension $2$, while at other
points it has dimension $1$ (where $x \neq 0$ or $y \neq 0$, the relation is nonzero modulo the maximal ideal, in characteristic $\neq 2, 3$). As $X$ is
reduced, Proposition~\ref{schemes:proposition-fibre-dimension}(3) shows that $\Omega$ is not locally free.
\end{solution}

\subsection{Cohomology of Coherent Sheaves}\label{schemes:subsection-exercises-coherent-cohomology}

\begin{exercise}\label{schemes:exercise-punctured-plane}
Compute $H^1(\Aa^2_k \setminus \{0\}, \mathcal{O})$ with the \v{C}ech complex of the cover $D(x) \cup D(y)$, and deduce again that $\Aa^2_k \setminus \{0\}$ is not affine.
\end{exercise}

\begin{solution}
The complex is $k[x, y]_x \times k[x, y]_y \to k[x, y]_{xy}$, $(a, b) \mapsto b - a$. Its cokernel has basis the monomials $x^{-i}y^{-j}$ with $i, j \geq 1$, so $H^1$ is
infinite-dimensional. The scheme is quasi-compact and separated, so the \v{C}ech complex computes cohomology
(Corollary~\ref{schemes:corollary-cech-computes}), and it is not affine by Theorem~\ref{schemes:theorem-affine-vanishing}.
\end{solution}

\begin{exercise}\label{schemes:exercise-plane-curve-cohomology}
Let $C \subseteq \PP^2_k$ be a plane curve of degree $d \geq 1$. Show that $h^0(\mathcal{O}_C) = 1$ and $h^1(\mathcal{O}_C) = \frac{(d-1)(d-2)}{2}$.
\end{exercise}

\begin{solution}
The sequence $0 \to \mathcal{O}(-d) \to \mathcal{O} \to \mathcal{O}_C \to 0$ and Theorem~\ref{schemes:theorem-projective-space} give
$0 \to k \to H^0(\mathcal{O}_C) \to H^1(\mathcal{O}(-d)) = 0 \to 0$ and $0 \to H^1(\mathcal{O}_C) \to H^2(\mathcal{O}(-d)) \to H^2(\mathcal{O}) = 0$. The space $H^2(\PP^2, \mathcal{O}(-d))$ has
basis the monomials $x^ay^bz^c$ with $a, b, c \leq -1$ and $a + b + c = -d$, and there are $\binom{d-1}{2}$ of them.
\end{solution}

\begin{exercise}\label{schemes:exercise-euler-projective}
Show that $\chi(\PP^r, \mathcal{O}(n)) = \binom{n+r}{r}$ for all $n \in \ZZ$, interpreting the binomial coefficient as the polynomial $\frac{(n+1)\cdots(n+r)}{r!}$, and check the
value at $n = -r-1$ against Serre duality.
\end{exercise}

\begin{solution}
For $n \geq 0$ only $H^0$ contributes and has dimension $\binom{n+r}{r}$. For $-r-1 < n < 0$ all cohomology vanishes and the polynomial vanishes too. For $n \leq -r-1$, only $H^r$
contributes, with dimension equal to the number of monomials with exponents $\leq -1$ summing to $n$, which is $\binom{-n-1}{r}$. This equals $(-1)^r\binom{n+r}{r}$
as polynomials in $n$. At $n = -r-1$ we get $(-1)^r$, matching $h^r(\omega) = h^0(\mathcal{O}) = 1$.
\end{solution}

\begin{exercise}\label{schemes:exercise-degree-hyperplane}
Let $X \subseteq \PP^r_k$ be a projective variety of dimension $d \geq 1$ over an infinite field, and $H$ a hyperplane not containing $X$. Show that $X \cap H$ has the same
degree as $X$.
\end{exercise}

\begin{solution}
From $0 \to \mathcal{O}_X(n - 1) \to \mathcal{O}_X(n) \to \mathcal{O}_{X \cap H}(n) \to 0$ we get $P_{X \cap H}(n) = P_X(n) - P_X(n-1)$. If $P_X$ has leading term $e\,n^d/d!$, the
difference has leading term $e\,n^{d-1}/(d-1)!$. So both have degree $e$.
\end{solution}
